\documentclass[11pt]{article}
\pdfoutput=1
\usepackage{jcapmod}
\usepackage{lmodern}
\usepackage{booktabs}
\usepackage{longtable}
\usepackage{bm}
\usepackage[english]{babel}
\usepackage[makeroom]{cancel}
\usepackage[per-mode=reciprocal, mode=math, group-digits=integer]{siunitx}[=v2]
\sisetup{range-units=brackets, range-phrase={,}, open-bracket=[, close-bracket=]}
\usepackage{subfig}
\usepackage[babel=true]{microtype}

\setcounter{tocdepth}{3}
\setlength{\textwidth}{460pt}
\setlength{\topmargin}{-1.2cm} \setlength{\textheight}{640pt} \setlength{\oddsidemargin}{10pt} \linespread{1.1}
\setlength{\parindent}{0.2in}

\numberwithin{equation}{section}
\allowdisplaybreaks[1]
\toccontinuoustrue
\newcommand{\GeV}{\,\text{GeV}}
\newcommand{\TeV}{\,\text{TeV}}
\newcommand{\MeV}{\,\text{MeV}}
\newcommand{\cm}{\,\text{cm}}
\newcommand{\Mpl}{M_{\text{Pl}}}
\newcommand{\Och}{\Omega_c h^2}
\newcommand{\sigmav}{\langle\sigma v\rangle}
\newcommand{\sigmavv}{\langle\sigma v^2\rangle}
\newcommand{\gstar}{g_*}
\newcommand{\gstarS}{g_{*S}}
\newcommand{\gtilde}{\tilde{g}}
\newcommand{\TD}{T_h}
\newcommand{\TSM}{T_{\text{SM}}}
\newcommand{\mX}{m_X}
\newcommand{\mY}{m_Y}
\newcommand{\alphaX}{\alpha_X}
\newcommand{\alphaEM}{\alpha_{\text{EM}}}
\newcommand{\epsX}{\varepsilon}
\newcommand{\Yeq}{Y^{\text{eq}}}
\newcommand{\neqMB}{n^{\text{eq}}_{\rm MB}}
\newcommand{\nequ}{n^{\text{eq}}}

\title{%
Documentation of the Hidden Sector Dark Matter Toolkit\\[0.5em]
\large Code Structure, Boltzmann Equations, and Assumptions
}
\author{Gabriele Montefalcone}
\date{\today}

\begin{document}
\maketitle

%======================================================================
\section{Introduction}
\label{sec:intro}
%======================================================================

This document provides a detailed technical description of the \texttt{secluded\_DM} Python toolkit, which computes the relic abundance of dark matter (DM) in \emph{hidden sector} scenarios. In these models, the DM particle~$X$ annihilates primarily into unstable mediator particles~$Y$, which subsequently decay to Standard Model (SM) final states through a small portal coupling~$\epsX$. The defining feature is the parametric hierarchy
\begin{equation}
    \alpha_X \gg \epsX^2\,,
\end{equation}
so that the DM annihilation rate $X\bar{X}\to YY$ is set by the dark fine-structure constant~$\alphaX$, while the mediator decay rate $Y\to \text{SM}$ is controlled by the much smaller portal coupling.

The code solves the coupled Boltzmann equations for the DM yield~$Y_X$, the mediator yield~$Y_Y$, and the dark-to-SM temperature ratio $\xi = \TD / \TSM$, incorporating:
\begin{itemize}
    \item $2\to2$ annihilation: $X\bar{X} \leftrightarrow YY$
    \item $3\to2$ cannibal processes: $YYY\to YY$, $YYX\to YX$, $YXX\to XX$
    \item Mediator decay: $Y\to\text{SM}$ with portal coupling $\epsX$
    \item Hidden sector heating and entropy injection
\end{itemize}

Five portal models are implemented: kinetic mixing with hypercharge (\texttt{VectorPortal}), $B-L$ gauge boson mixing (\texttt{BLPortal}), $L_i - L_j$ gauge boson mixing (\texttt{LiLjPortal}), baryon-number gauge boson mixing (\texttt{BaryonPortal}), and scalar mixing with the SM Higgs boson (\texttt{HiggsPortal}). Observational constraints from the DM relic abundance, Big Bang Nucleosynthesis (BBN) hadronic injection bounds, and direct detection experiments (particularly the combination of LZ and XENONnT bounds) are computed for each model.


%======================================================================
\section{Cosmological Background}
\label{sec:cosmology}
%======================================================================
This section establishes the cosmological quantities that enter the Boltzmann equations of Sec.~\ref{sec:boltzmann}. We first define the Standard Model (SM) thermal bath and the Hubble rate, including the modification from the hidden-sector energy density, and then collect the equilibrium distributions and energy moments under the Maxwell--Boltzmann approximation that are used throughout the solver.
\subsection{Standard Model Thermal Bath and Hubble Rate}

The SM thermal background is described by the relativistic degrees of freedom $\gstar(T)$ and $\gstarS(T)$, loaded from tabulated data files (covering $T\sim 0.1\MeV$ to $\sim 100\GeV$) and interpolated linearly in $T$. The key thermodynamic quantities are:
\begin{align}
    \rho_{\text{rad}}(T) &= \frac{\pi^2}{30}\,\gstar(T)\,T^4 &
    &\text{(radiation energy density)} \label{eq:rhoR}\\
    s(T) &= \frac{2\pi^2}{45}\,\gstarS(T)\,T^3 &
    &\text{(entropy density)} \\
    H(T) &= \sqrt{\frac{\pi^2}{90}\,\gstar(T)}\;\frac{T^2}{\Mpl} &
    &\text{(Hubble rate, SM only)}
\end{align}
where $\Mpl = 2.435\times 10^{18}\GeV$ is the reduced Planck mass. The factor
\begin{equation}
    \gtilde(T) = 1 + \frac{1}{3}\frac{d\ln \gstarS}{d\ln T}\label{eq:gtilde}
\end{equation}
accounts for the temperature dependence of $\gstarS$ in the entropy conservation equation ($\gtilde = 1$ when $\gstarS$ is constant).

When the hidden sector contributes non-negligibly to the total energy density, the Hubble rate is modified:
\begin{equation}
    H_{\text{tot}} = \sqrt{H^2 + \frac{\rho_h}{3\Mpl^2}}\,, \qquad
    \rho_{h} = \rho_X +\rho_Y. \label{eq:Htot}
\end{equation}

\subsection{Equilibrium Distributions and Energy Moments}
\label{sec:bessel}

All equilibrium quantities use the Maxwell-Boltzmann approximation (see Sec.~\ref{sec:assumptions} for a discussion of its validity). For a species with mass $m$, internal degrees of freedom $g$, at a temperature $T$, the number density and its derivative are given by:
\begin{equation}
    \neqMB(z) = \frac{g\,m^3}{2\pi^2}\,\frac{K_2(z)}{z}\,, \qquad
    \frac{d\neqMB}{dT} = \frac{g\,m^2}{2\pi^2}\Big[z\,K_1(z) + 3\,K_2(z)\Big]
\end{equation}
where $z\equiv m/T$ and $K_\nu(z)$ is the modified Bessel functions of the second kind. A useful combination is the logarithmic derivative
\begin{equation}
    \lambda(z) \equiv \frac{d\ln \neqMB}{d\ln T} = z\,\frac{K_1(z)}{K_2(z)} + 3\,,\label{eq:lambda}
\end{equation}
which appears in the Boltzmann equations whenever equilibrium tracking is needed.

Higher moments of the distribution give the energy density and pressure,
\begin{equation}
    \rho = B_1\,\neqMB \, , \qquad \rho + P = B_2\,\neqMB \, , \label{eq:MB_rho_P}
\end{equation}
where
\begin{align}
    B_1(z, m) &= m\!\left(\frac{K_1(z)}{K_2(z)} + \frac{3}{z}\right), \qquad
    B_2(z, m) = m\!\left(\frac{K_1(z)}{K_2(z)} + \frac{4}{z}\right) \label{eq:B1}
\end{align}
which can be thought of as the average energy per particle, $B_1 = \langle E\rangle / n$, and the enthalpy per particle, $B_2 = (\rho + P)/n$, respectively. Note that the difference $B_2 - B_1 = m/z = T$ recovers the ideal-gas relation $P = nT$. Their temperature derivatives are
\begin{align}
    \frac{dB_1}{dT} &= 3 - z^2\!\left[R^2 - 1 + \frac{3R}{z}\right], \label{eq:dB1dT} \\
    \frac{dB_2}{dT} &= 4 - z^2\!\left[R^2 - 1 + \frac{3R}{z}\right],\label{eq:dB2dT}
\end{align}
where $R \equiv K_1(z)/K_2(z)$.

Another useful quantity that is implemented in the Boltzmann equations is the logarithmic equilibrium yield, $\ln Y_i^{\text{eq}} = \ln \neqMB - \ln s$, with $\ln 2$ added if antiparticles are included.

\paragraph{Code implementation: asymptotic expansions.}
In our code, the density is set to zero for $z > 500$ to avoid underflow and is doubled when antiparticles are present. For $z \geq 300$, the equilibrium number density is computed via the asymptotic form
\begin{equation}
    \ln \neqMB = \ln\!\Big(\frac{g\,m^3}{2\pi^2}\Big) + \frac{1}{2}\ln\frac{\pi}{2z} - z + \ln\!\Big(1 + \frac{15}{8z} + \frac{105}{128z^2}\Big) - \ln z\,.
\end{equation}
The Bessel ratio $R = K_1(z)/K_2(z)$, which enters $B_1$, $B_2$, $\lambda$, and their derivatives, is computed directly from \texttt{scipy.special.kv} for $z < 50$ while for $z \geq 50$ the following expansion is used instead:
\begin{equation}
    \frac{K_1(z)}{K_2(z)} \approx 1 - \frac{3}{2z} + \frac{15}{8z^2} - \frac{105}{16z^3}.
\end{equation}
In the same regime, the bracket in the $B$-factor derivatives simplifies to $3/(2z^2) - 15/(4z^3)$.
 
\paragraph{Full quantum statistics (\texttt{species\_thermo}):} A separate method computes $(\rho, P, n)$ with full Fermi-Dirac or Bose-Einstein integrals using a three-regime approach (ultra-relativistic analytic for $z \leq 10^{-3}$, numerical integral for $10^{-3} < z < 1$, MB for $z \geq 1$). This is available for auxiliary calculations but is \emph{not} used in the Boltzmann solver itself.


%======================================================================
\section{Hidden Sector Models}
\label{sec:models}
%======================================================================

This section presents the five portal models implemented in the code. Each model specifies the hidden-sector particle content (masses $\mX$ and $\mY$, internal degrees of freedom $g_X$
$g_Y$, and whether antiparticles are included), the $2\to2$  annihilation and $3\to2$ cannibal cross sections that govern the freeze-out dynamics, the mediator decay width $\Gamma(Y\to\mathrm{SM})$, and the spin-independent DM--nucleon scattering cross section $\sigma_{\rm SI}$. All models inherit from the abstract base class \texttt{DarkSectorModel}, which defines this common interface, and differ only in how the portal coupling $\epsX$ connects the hidden sector to the Standard Model. The initial dark-to-SM temperature ratio $\xi_{\rm ini}$ is a free parameter common to all models.

\subsection{\texttt{VectorPortal}: Kinetic Mixing with Hypercharge}
\label{sec:VP}

The DM particle $X$ is a Dirac fermion ($g_X = 2$, includes antiparticles) charged under a dark $U(1)_D$ gauge symmetry. The mediator $Y \equiv Z'$ is the associated dark photon ($g_Y = 3$, massive spin-1, no antiparticle doubling). The dark photon mixes kinetically with the SM hypercharge gauge boson via the mixing parameter~$\epsX$. We first describe the hidden-sector processes --- $2\to2$ annihilations and $3\to2$ cannibal reactions --- and then the connection to the visible sector through kinetic mixing, which controls both the mediator decay width and the DM--nucleon scattering cross section.

\subsubsection{Hidden-sector processes}
\label{sec:VP_dark}

The thermally averaged $2\to2$ annihilation cross section for $X\bar{X}\to Z'Z'$ is
\begin{equation}
    \sigmav_{X\bar{X}\to Z'Z'} = 2\!\left(a + \frac{6b}{x_X}\right),
\end{equation}
where $x_X = \mX/\TD$, $r_v \equiv \mY/\mX$, and
\begin{align}
    a &= \frac{2\pi\alphaX^2}{\mX^2}\,\frac{(1 - r_v^2)^{3/2}}{(2 - r_v^2)^2}\,, \label{eq:VP_swave} \\
    b &= \frac{\pi\alphaX^2}{12\,\mX^2}\,\frac{\sqrt{1 - r_v^2}\,(24 + 28r_v^2 - 36r_v^4 + 17r_v^6)}{(2 - r_v^2)^4} \label{eq:VP_pwave}
\end{align}
are respectively the $s$-wave and $p$-wave coefficients. For the background evolution solver, where the $p$-wave contribution is subdominant at late times ($x \gg 1$), only the $s$-wave piece $\sigmav_{\text{s-wave}} = 2a$ is retained.

Three $3\to2$ cannibal processes contribute to number-changing interactions within the hidden sector. The dominant channel for the freeze-out of cannibal processes is $YYX\to YX$, whose thermally averaged cross section
\begin{equation}
    \sigmavv_{YYX\to YX} = \Delta_3\,\frac{\alphaX^3}{\mX^5}
    \label{eq:sv2_YYX}
\end{equation}
is computed exactly at tree level, with the coefficient
\begin{equation}
    \Delta_3 = \frac{4\pi^2\,r^5\,\sqrt{\tfrac{4}{3}r(r+2)+1}\;\mathcal{N}_3(r)}{(r+2)^3\,(2r+1)^4\,\big[2r^2(r+2) + r - 1\big]^2}
    \label{eq:Delta3}
\end{equation}
where $r = \mX/\mY$ and
\begin{equation}
    \mathcal{N}_3(r) = 2r\Big(r\big(r(r(2r(r(8r(10r+71)+1683)+2798)+6107)+4722)+2335\big)+578\Big)+195\,.
\end{equation}
The remaining two channels are~\cite{Berlin:2016gtr}:
\begin{align}
    \sigmavv_{YYY\to YY} &= \Delta_1\,\frac{\alphaX^5\,\TD^7}{\mX^{12}}\,, \label{eq:sv2_YYY} \\
    \sigmavv_{YXX\to XX} &= \Delta_2\,\frac{\alphaX^3}{\mX^5}\,, \label{eq:sv2_YXX}
\end{align}
where $\Delta_2$ is also computed exactly at tree level:
\begin{equation}
    \Delta_2 = \frac{\pi^2\,r^3\,(1+4r)^{3/2}\;\mathcal{N}_2(r)}{6\,(1+2r)^3\,(2r^2 + r - 1)^2}
    \label{eq:Delta2}
\end{equation}
with
\begin{equation}
    \mathcal{N}_2(r) = 256r^6 - 384r^5 + 536r^4 - 96r^3 + 38r^2 + 72r + 31\,.
\end{equation}
The $YYY\to YY$ rate is loop-suppressed relative to the other two channels ($\alphaX^5$ vs.\ $\alphaX^3$) and carries an explicit temperature dependence, making it negligible for the evolution of the hidden sector. We retain it for completeness with a phenomenological coefficient $\Delta_1 = 0.5$.

\subsubsection{Connection to the Standard Model}
\label{sec:VP_SM}

After restoring canonical kinetic terms, the $Z'$ acquires couplings to SM fermions proportional to $\epsX$. For a fermion $f$ with electric charge $Q_f$ and weak hypercharges $Y_L$, $Y_R$, the left- and right-handed couplings are~\cite{Berlin:2016gtr}
\begin{align}
    g_L &= \epsX\,\frac{m_{Z'}^2\,g_Y^{\text{SM}}\,Y_L - M_Z^2\,g_W\,s_W\,c_W\,Q_f}{M_Z^2 - m_{Z'}^2}\,, \label{eq:gL} \\
    g_R &= \epsX\,\frac{m_{Z'}^2\,g_Y^{\text{SM}}\,Y_R - M_Z^2\,g_W\,s_W\,c_W\,Q_f}{M_Z^2 - m_{Z'}^2}\,, \label{eq:gR}
\end{align}
where $g_W$ and $g_Y^{\text{SM}}$ are the $SU(2)_L$ and $U(1)_Y$ gauge couplings, $s_W \equiv \sin\theta_W$ and $c_W \equiv \cos\theta_W$ with $\theta_W$ the Weinberg angle, and $M_Z$ is the $Z$-boson mass. The vector and axial couplings to each fermion species are then
\begin{equation}
    g_{fV} = \frac{g_R + g_L}{2}\,, \qquad g_{fA} = \frac{g_R - g_L}{2}\,.
\end{equation}

\paragraph{Mediator decay width.}
These couplings determine the $Z'$ decay width. The partial width into a fermion pair is
\begin{equation}
    \Gamma(Z'\to f\bar{f}) = \frac{N_c^f\,m_{Z'}}{12\pi}\,\sqrt{1 - r_f^2}\;\Big[g_{fV}^2\,\big(1 + \tfrac{1}{2}r_f^2\big) + g_{fA}^2\,\big(1 - r_f^2\big)\Big]
    \label{eq:VP_decay}
\end{equation}
where $r_f = 2m_f/m_{Z'}$ and $N_c^f$ is the colour factor ($3$ for quarks, $1$ for leptons). The total width is the sum over all kinematically accessible channels ($m_{Z'} > 2m_f$):
\begin{equation}
    \Gamma_{Z'}^{\text{tot}} = \sum_{f:\,m_{Z'}>2m_f} \Gamma(Z'\to f\bar{f})\,.
\end{equation}

\paragraph{Direct-detection cross section.}
The same kinetic mixing generates a spin-independent DM--nucleon scattering cross section through $Z$--$Z'$ exchange, following~\cite{Hu:2025thq}. The $Z$--$Z'$ mixing angle $\eta$ is defined by
\begin{equation}
    \tan(2\eta) = \frac{-2\epsX\,s_W\,M_Z^2}{M_Z^2 - m_{Z'}^2}\,,
\end{equation}
and the DM couplings to the mass eigenstates are $V_\chi^Z = g_D\sin\eta$ and $V_\chi^{Z'} = g_D\cos\eta$, with $g_D = \sqrt{4\pi\alphaX}$. The quark couplings receive contributions from both the electromagnetic and weak neutral currents:
\begin{align}
    V_q^Z &= \sqrt{4\pi\alphaEM}\,\epsX\,\sin\eta\,Q_q + \frac{\sqrt{4\pi\alphaEM}}{s_W c_W}\,(\cos\eta + \epsX\,s_W\,\sin\eta)\,(T_3^q - s_W^2\,Q_q)\,, \\
    V_q^{Z'} &= \sqrt{4\pi\alphaEM}\,\epsX\,\cos\eta\,Q_q + \frac{\sqrt{4\pi\alphaEM}}{s_W c_W}\,(\sin\eta - \epsX\,s_W\,\cos\eta)\,(T_3^q - s_W^2\,Q_q)\,,
\end{align}
where $\alphaEM \simeq 1/137$ is the electromagnetic fine-structure constant, $Q_q$ the electric charge of quark $q$, and $T_3^q$ its weak isospin. The spin-independent cross section per nucleon is then
\begin{equation}
    \sigma_{\text{SI}} = \frac{\mu^2}{\pi}\left[\sum_{i\in\{Z,Z'\}}\frac{V_\chi^i\;\big[\mathcal{Z}(2V_u^i + V_d^i) + (\mathcal{A}-\mathcal{Z})(V_u^i + 2V_d^i)\big]}{\mathcal{A}\,m_i^2}\right]^2 \times \hbar^2 c^2
    \label{eq:VP_sigma_SI}
\end{equation}
where $\mu = \mX\,m_N/(\mX + m_N)$ is the reduced mass with $m_N = (m_p + m_n)/2$ the average nucleon mass, and $\mathcal{Z}$ and $\mathcal{A}$ are the atomic and mass numbers of the target nucleus.

Two notable features appear in the cross section as a function of $m_{Z'}$. First, the couplings in Eqs.~\eqref{eq:gL}--\eqref{eq:gR} exhibit a resonance at $m_{Z'} = M_Z$, where $\sigma_{\text{SI}}$ is strongly enhanced and the required $\epsX$ is correspondingly minimised. In practice this pole would be regularised by the finite width $\Gamma_{Z'}$, which we neglect here; nevertheless, the enhancement remains significant in the vicinity of the $Z$ pole. Second, the coherent sum over $Z$ and $Z'$ exchange leads to an exact cancellation between the two amplitudes above the pole. For $m_{Z'} > M_Z$ the mixing angle $\eta$ changes sign, flipping the relative sign of the two amplitudes. At a critical mass
\begin{equation}
    m_{Z'}^{\cancel{\sigma}} \simeq 129.4~\text{GeV}
    \;\;\bigl(\approx 1.42\,M_Z\bigr)
    \qquad\text{(Xe-131 target)},
    \label{eq:VP_blind_spot}
\end{equation}
the two amplitudes cancel exactly and $\sigma_{\text{SI}} = 0$, independent of $\epsX$ and $\alphaX$ to leading order. At this blind spot the required $\epsX \to \infty$, making the model undetectable. The position of the zero depends on the target nucleus through the $(\mathcal{Z}, \mathcal{A})$-weighted quark couplings but is independent of the hidden-sector parameters.

The default parameter values used in the \texttt{VectorPortal} class are summarised in Table~\ref{tab:VP_defaults}.
\begin{table}[h]
\centering
\begin{tabular}{lll}
\toprule
Parameter & Symbol & Default \\
\midrule
DM mass & $\mX$ & $100\GeV$ \\
Mediator mass & $\mY$ & $30\GeV$ \\
Dark coupling & $\alphaX$ & $10^{-2}$ \\
DM DOF & $g_X$ & 2 (Dirac fermion) \\
Mediator DOF & $g_Y$ & 3 (massive vector) \\
Cannibal rescaling & $\Delta_1$ & 0.5 \\
Cannibal rescaling & $\Delta_2, \Delta_3$ & exact tree-level \\
Initial temperature ratio & $\xi_{\text{ini}}$ & 1.0 \\
\bottomrule
\end{tabular}
\caption{Default parameters for the \texttt{VectorPortal} model.}
\label{tab:VP_defaults}
\end{table}
\subsection{\texttt{BLPortal}: $B-L$ Kinetic Mixing}
\label{sec:BL}

The hidden-sector $Z'$ mixes kinetically with the $U(1)_{B-L}$ gauge boson $Z_{B-L}$. The coupling to SM fermion $f$ is purely vectorial~\cite{Hooper:2019xss}:
\begin{equation}
    g_{fV} = \epsX\,g_{B-L}\,(B-L)_f\;\left|\frac{m_{Z_{B-L}}^2 + m_{Z'}^2}{m_{Z_{B-L}}^2 - m_{Z'}^2}\right|, \qquad g_{fA} = 0
    \label{eq:BL_coupling}
\end{equation}
where $g_{B-L}$ is the $B-L$ gauge coupling and $(B-L)_f$ the fermion charge: $+1/3$ for quarks and $-1$ for leptons. Since the hidden sector is identical to the \texttt{VectorPortal}, all $2\to2$ and $3\to2$ cross sections (Eqs.~\eqref{eq:VP_swave}--\eqref{eq:sv2_YYX}) carry over unchanged. The $Z'$ decay width retains the same functional form as Eq.~\eqref{eq:VP_decay}, but with purely vectorial couplings ($g_{fA} = 0$) and all kinematically accessible fermions contributing. For direct detection, all quarks carry the same $(B-L) = 1/3$ charge, so the coupling is isospin-universal: proton and neutron receive equal contributions. Evaluating $g_{qV}$ from Eq.~\eqref{eq:BL_coupling} with $(B-L)_q = 1/3$, the three valence quarks in a nucleon contribute coherently, giving
\begin{equation}
    \sigma_{\text{SI}} = \frac{\mu^2}{\pi}\left(\frac{3\,g_D\,g_{qV}}{m_{Z'}^2}\right)^2 \times \hbar^2 c^2
    \label{eq:BL_sigma_SI}
\end{equation}
where $g_D$ and $\mu$ are as defined in Sec.~\ref{sec:VP_SM}. Unlike the hypercharge portal, only a single mediator contributes and no blind spot arises.

\subsection{\texttt{BaryonPortal}: Baryon Number Kinetic Mixing}
\label{sec:Baryon}

The $Z'$ mixes kinetically with the gauge boson of $U(1)_B$, acquiring purely vectorial couplings to quarks proportional to their baryon number~\cite{Hooper:2019xss}:
\begin{equation}
    g_{fV} = \epsX\,g_B\,B_f\;\left|\frac{m_{Z_B}^2 + m_{Z'}^2}{m_{Z_B}^2 - m_{Z'}^2}\right|, \qquad g_{fA} = 0
    \label{eq:B_coupling}
\end{equation}
where $g_B$ is the $U(1)_B$ gauge coupling, $m_{Z_B}$ the associated gauge boson mass, and $B_f = 1/3$ for all quarks. Leptons are uncharged, so the $Z'$ decays exclusively to hadronic channels. The hidden-sector cross sections are again identical to the \texttt{VectorPortal}. Since all quarks carry the same baryon charge, the direct-detection cross section is isospin-universal and takes the same form as Eq.~\eqref{eq:BL_sigma_SI}, with $g_{qV}$ evaluated from Eq.~\eqref{eq:B_coupling}.

\subsection{\texttt{LiLjPortal}: $L_i - L_j$ Kinetic Mixing}
\label{sec:LiLj}

The $Z'$ mixes kinetically with the gauge boson of $U(1)_{L_i - L_j}$, acquiring purely vectorial couplings to leptons and their neutrinos of the two involved families~\cite{Hooper:2019xss}: flavour $i$ carries charge $+1$ and flavour $j$ carries charge $-1$, while quarks are uncharged. Three flavour combinations are supported: $e$--$\mu$, $\mu$--$\tau$, and $e$--$\tau$. The hidden-sector cross sections are identical to those of the \texttt{VectorPortal} (Eqs.~\eqref{eq:VP_swave}--\eqref{eq:sv2_YYX}). Because quarks are uncharged, the $Z'$ decays exclusively to leptons at tree level and the spin-independent DM--nucleon cross section vanishes, $\sigma_{\text{SI}} = 0$. A nonzero $\sigma_{\text{SI}}$ arises at loop level, but the resulting cross section is too small to produce meaningful direct-detection constraints unless the couplings are very large, so we neglect it.

\subsection{\texttt{HiggsPortal}: Scalar Mixing}
\label{sec:HP}

The DM particle $X$ is a Majorana fermion ($g_X = 2$, no antiparticle doubling) coupled to a real scalar mediator $\phi$ ($g_Y = 1$) that mixes with the SM Higgs boson through the mixing angle $\epsX$, following the conventions of~\cite{Hooper:2019xss}. In the small-angle limit $\sin\epsX \approx \epsX$. Two couplings parametrise the hidden-sector Yukawa interactions:
\begin{equation}
    \mathcal{L} \supset -\frac{\lambda_s}{2}\,\phi\,\bar{X}X - \frac{i\lambda_p}{2}\,\phi\,\bar{X}\gamma_5 X
\end{equation}
where $\lambda_s$ and $\lambda_p$ are the scalar and pseudoscalar couplings respectively. We first describe the hidden-sector processes and then the connection to the SM through scalar mixing, which controls both the mediator decay width and the DM--nucleon scattering cross section.

\subsubsection{Hidden-sector processes}
\label{sec:HP_dark}

The thermally averaged $2\to2$ annihilation cross section for $X\bar{X}\to\phi\phi$ is
\begin{equation}
    \sigmav_{X\bar{X}\to\phi\phi} = \!\left(a + \frac{6b}{x_X}\right),
\end{equation}
where $x_X = \mX/\TD$, $r \equiv \mY/\mX$, and
\begin{align}
    a &= \frac{2\sqrt{1-r^2}\,\lambda_p^2\,\lambda_s^2}{\mX^2\,\pi\,(r^2 - 2)^2}\,, \label{eq:HP_swave} \\
    b &= \frac{1}{12\,\mX^2\,\pi\,\sqrt{1-r^2}\,(r^2-2)^4}\Big[
    -2(r^2-1)^3\,\lambda_p^4
    + 3(r^6 - 8r^4 + 20r^2 - 12)\,\lambda_s^2\,\lambda_p^2 \notag \\
    &\qquad + 2(-2r^6 + 10r^4 - 17r^2 + 9)\,\lambda_s^4\Big] \label{eq:HP_pwave}
\end{align}
are respectively the $s$-wave and $p$-wave coefficients. Note that the $s$-wave vanishes when either $\lambda_s = 0$ or $\lambda_p = 0$, so both couplings must be nonzero for a nonvanishing leading-order cross section. For the background evolution solver, only the $s$-wave piece $\sigmav_{\text{s-wave}} = a$ is retained, since at late times ($x \gg 1$) the $p$-wave contribution is Boltzmann-suppressed and negligible. For $x_X < 1$ the $p$-wave term is also dropped, which is inconsequential since the hidden sector is still in full equilibrium in this regime.

Three $3\to2$ cannibal processes are included, all computed exactly at tree level. In analogy to the \texttt{VectorPortal} (Eq.~\eqref{eq:sv2_YYX}), the dominant channel for cannibal freeze-out is $\phi\phi X\to\phi X$, whose cross section can be written as
\begin{equation}
    \sigmavv_{\phi\phi X\to\phi X} = \Delta_3(r,\rho)\,\frac{\lambda_s^3\,\lambda_p^3}{\mX^5}
    \label{eq:HP_sv2_YYX}
\end{equation}
where $r \equiv \mX/\mY$, $\rho \equiv \lambda_p/\lambda_s$, and
\begin{multline}
    \Delta_3(r,\rho) = \frac{9\sqrt{3}\,r^5\,\sqrt{3+8r+4r^2}}{64\pi\,(r+2)^3\,(2r+1)^4\,(2r^3+4r^2+r-1)^2}
    \times \Big[
    (1+2r)^5(3+2r) \\
    + (9 + 16r + 180r^2 + 400r^3 + 288r^4 + 64r^5)(1+2r)^2\rho^2
    + 3(-1+2r+2r^2)^2\rho^6 \\
    + (9 + 8r + 140r^2 + 888r^3 + 2220r^4
    + 3360r^5 + 3024r^6 + 1408r^7 + 256r^8)\rho^4
    \Big]\,.
    \label{eq:HP_Delta3}
\end{multline}
The second channel, $\phi XX\to XX$, shares the same factorisation:
\begin{equation}
    \sigmavv_{\phi XX\to XX} = \Delta_2(r,\rho)\,\frac{\lambda_s^3\,\lambda_p^3}{\mX^5}
    \label{eq:HP_sv2_YXX}
\end{equation}
with
\begin{multline}
    \Delta_2(r,\rho) = \frac{r^3\,\sqrt{1+4r}}{64\pi\,(1+2r)^3\,(-1+r+2r^2)^2}
    \times \Big[
    (1+8r)(1+\rho^2)^3
    + 4r^2(2 + 21\rho^2 + 20\rho^4 + \rho^6) \\
    - 16r^3(4 - 7\rho^2 - 6\rho^4 + 5\rho^6)
    - 16r^4(7 - 6\rho^2 - 25\rho^4 + 4\rho^6) \\
    + 128r^5(1 + 9\rho^2 + 14\rho^4 + 2\rho^6)
    + 256r^6(1 + 2\rho^2)^2
    \Big]\,.
    \label{eq:HP_Delta2}
\end{multline}
This channel is generally subdominant but is included for completeness. The third channel, $\phi\phi\phi\to\phi\phi$, factorises with a different coupling dependence:
\begin{equation}
    \sigmavv_{\phi\phi\phi\to\phi\phi} = \Delta_1(r,\rho)\,\frac{\lambda_s^5\,\lambda_p^5}{\mX^5}
    \label{eq:HP_sv2_YYY}
\end{equation}
with
\begin{equation}
    \Delta_1(r,\rho) = \frac{\sqrt{5}}{3\pi^5}\,\frac{r^3\,(3 + 10\rho^2 + 15\rho^4)^2}{\rho^3}\,.
    \label{eq:HP_Delta1}
\end{equation}
Unlike the tree-level channels above, $\phi\phi\phi\to\phi\phi$ proceeds through a one-loop pentagon diagram (a closed $X$ loop with five external $\phi$ legs) and the expression above is derived in the heavy-fermion limit $r \gg 1$, where all external momenta are negligible compared to the loop momentum. 
This is self-consistent as the $\phi\phi\phi\to\phi\phi$ channel can only become competitive with $\phi\phi X\to\phi X$ precisely in this regime, since at low $r$ the dark matter $X$ is comparably abundant to $\phi$ and the tree-level process wins without any loop penalty. At large $r$, however, $X$ is exponentially Boltzmann-suppressed in the initial state, and the cannibal process --- which only requires $X$ as a virtual particle inside the loop --- sidesteps this suppression, compensating the $\sim\!\lambda^4/(16\pi^2)^2$ loop cost.

\subsubsection{Connection to the Standard Model}
\label{sec:HP_SM}

\paragraph{Mediator decay width.}
The mediator $\phi$ decays to SM particles through its mixing with the Higgs boson. All partial widths are proportional to $\sin^2\!\epsX \approx \epsX^2$, and we define the normalised total width $\Gamma_{\text{norm}}(m_\phi) \equiv \Gamma_{\text{tot}}/\sin^2\!\epsX$ so that the physical width is
\begin{equation}
    \Gamma(\phi\to\text{SM}) = \epsX^2\;\Gamma_{\text{norm}}(m_\phi)\,.
\end{equation}
The SM-like contribution to $\Gamma_{\text{norm}}$ is computed using a hybrid approach across four mass regimes:
\begin{itemize}
    \item $m_\phi < 5\GeV$: the \texttt{scalar\_portal} package~\cite{Winkler:2018qyg}, which uses a dispersive analysis below $2\GeV$ (leptonic decays, hadronic channels via form factors, multi-meson final states) and perturbative QCD with running quark masses ($N_f = 4$) above $2\GeV$.
    \item $5\GeV \leq m_\phi \leq 500\GeV$: the \texttt{HDECAY} Fortran program~\cite{Djouadi:1997yw, Djouadi:2018xqq}, which includes higher-order QCD corrections, off-shell decays to $WW^*/ZZ^*$, and loop-induced channels ($gg$, $\gamma\gamma$).
    \item $m_\phi > 500\GeV$: tree-level decay widths with exact kinematics, since \texttt{HDECAY} is not reliable in this regime due to un-resummed Sudakov logarithms.
\end{itemize} If $m_\phi > 2m_h$, an additional contribution from $\phi\to hh$ arises in the singlet extension:
\begin{equation}
    \frac{\Gamma(\phi\to hh)}{\sin^2\!\epsX} = \frac{(2m_\phi^2 + m_h^2)^2}{128\pi\,m_\phi\,v_H^2}\,\sqrt{1 - \frac{4m_h^2}{m_\phi^2}}
\end{equation}
with $v_H = 246.22\GeV$ the Higgs vacuum expectation value. This channel can be switched off by setting \texttt{hh\_coupling=`zero'}.

\paragraph{Direct-detection cross section.}
The spin-independent DM--nucleon scattering cross section proceeds via $t$-channel exchange of the two mixed scalars $h$ and $\phi$. After electroweak symmetry breaking, the mass mixing $\phi \to \phi - \epsX\,h$ generates an effective $\phi$--nucleon coupling proportional to $\epsX\,\lambda_s$. In the small-mixing limit, the $h$- and $\phi$-exchange amplitudes add coherently with opposite sign (since the mixing rotates them in opposite directions), giving a per-nucleon amplitude proportional to the propagator difference $1/\mY^2 - 1/m_h^2$. The resulting cross section is
\begin{equation}
    \sigma_{\text{SI}} = \frac{4\mu^2}{\pi}\,a_N^2\,\hbar^2 c^2\,, \qquad
    a_N = \frac{\lambda_s\,\epsX\,f_N\,m_N}{v_H}\left(\frac{1}{\mY^2} - \frac{1}{m_h^2}\right)
    \label{eq:HP_sigma_SI}
\end{equation}
where $f_N \simeq 0.30$ is the Higgs--nucleon form factor, $m_h$ the SM Higgs mass, and $\mu$ the DM--nucleon reduced mass. The factor of 4 reflects the Majorana nature of $X$ (identical-particle combinatorics). This result is isospin-universal and independent of the target nucleus composition, since the coupling to nucleons proceeds entirely through the Higgs--nucleon vertex. The numerical estimate in~\cite{Berlin:2016gtr} (their Eq.~67) is consistent with this expression in the limit $\mY \ll m_h$. The propagator difference vanishes at $\mY = m_h$, where the two amplitudes cancel exactly and $\sigma_{\text{SI}} = 0$. At this blind spot the required $\epsX \to \infty$, making the model undetectable by direct-detection experiments regardless of the hidden-sector couplings.

The default parameter values used in the \texttt{HiggsPortal} class are summarised in Table~\ref{tab:HP_defaults}.
\begin{table}[h]
\centering
\begin{tabular}{lll}
\toprule
Parameter & Symbol & Default \\
\midrule
DM mass & $\mX$ & $100\GeV$ \\
Mediator mass & $\mY$ & $30\GeV$ \\
Dark couplings & $\lambda_s = \lambda_p$ & $10^{-2}$ \\
DM DOF & $g_X$ & 2 (Majorana fermion) \\
Mediator DOF & $g_Y$ & 1 (real scalar) \\
Cannibal rescaling & $\Delta_1, \Delta_2, \Delta_3$ & exact tree-level \\
Initial temperature ratio & $\xi_{\text{ini}}$ & 1.0 \\
$\phi\to hh$ coupling & \texttt{hh\_coupling} & \texttt{`singlet'} \\
\bottomrule
\end{tabular}
\caption{Default parameters for the \texttt{HiggsPortal} model.}
\label{tab:HP_defaults}
\end{table}

\section{Evolution of the Hidden Sector}\label{sec:evolution}
The cosmological evolution of the hidden sector can be naturally divided into two stages, each amenable to a distinct formulation.
The first stage tracks the coupled evolution of the DM particle $X$, the mediator $Y$, and the hidden-sector temperature as the hidden sector departs from internal chemical equilibrium. The key processes during this phase are the $2\to2$ annihilation $X\bar{X}\to YY$, which sets the DM relic abundance, and the $3\to2$ cannibal reactions, which maintain the mediator in chemical equilibrium until they too freeze out. Throughout this work, the portal coupling $\varepsilon$ is assumed to be sufficiently small that the hidden sector remains thermally decoupled from the Standard Model bath, evolving independently at its own temperature $T_h = \xi\,T_{\rm SM}$. In this regime, the mediator decay width generically satisfies $\Gamma_Y \ll H$ during the entire freeze-out process, so that it can be safely neglected and the dynamics is governed solely by the hidden-sector couplings. The evolution during this stage is then captured by a set of coupled Boltzmann equations, which we present in Sec.~\ref{sec:boltzmann}. We first write these equations in their standard, physically transparent form and then show how they can be recast in terms of chemical-potential variables and the hidden to SM temperature ratio $\xi$, a formulation that is better suited for numerical integration.

The second stage begins once the Boltzmann system has converged and the comoving abundances of $X$ and $Y$ are effectively frozen. The relevant dynamics is now the decay of the long-lived mediator into SM particles. As $Y$ depletes, it injects entropy into the visible sector and dilutes the dark-matter yield; it is through this mechanism that the portal coupling $\varepsilon$ ultimately controls the predicted relic abundance. This late-time background evolution is described in Sec.~\ref{sec:background}.

\subsection{Boltzmann Equations}
\label{sec:boltzmann}

The starting point is the standard set of Boltzmann equations for the
number densities $n_X$ and $n_Y$ and the total hidden-sector energy
density $\rho_h = \rho_X + \rho_Y$.  In a Friedmann--Robertson--Walker background these read:
\begin{align}
    &\dot{n}_X + 3H n_X = -\sigmav_{X}\left(n_X^2 - (n_X^{\text{eq}})^2\,\frac{n_Y^2}{(n_Y^{\text{eq}})^2}\right) \label{eq:boltz_nX} \\
    &\dot{n}_Y + 3H n_Y = \sigmav_{X}\left(n_X^2 - (n_X^{\text{eq}})^2\,\frac{n_Y^2}{(n_Y^{\text{eq}})^2}\right) - \Gamma_{\rm can}\,(n_Y - n_Y^{\text{eq}}) \label{eq:boltz_nY} \\
    &\dot{\rho}_h + 3H(\rho_h + P_h) = 0 \label{eq:boltz_rho}
\end{align}
where a dot corresponds to a derivative with respect to cosmic time, $H$ is the total Hubble rate including the hidden-sector contribution to the energy density, Eq.~\eqref{eq:Htot}, $\sigmav_{X}$ is the thermally averaged annihilation cross section, and $\Gamma_{\rm can}$ collects all $3\to2$ number-changing processes, namely:
\begin{equation}
  \Gamma_{\rm can} = \langle\sigma v^2\rangle_{YYY\to YY}\,n_Y^2\,+\langle\sigma v^2\rangle_{YYX\to YX}\,n_Y\,n_X+\,
\langle\sigma v^2\rangle_{YXX\to XX}\,n_X^2\,.
\label{eq:Gamma_can_std}
\end{equation}


The first equation governs the DM abundance, with the collision term vanishing when $n_X/n_Y = n_X^{\rm eq}/n_Y^{\rm eq}$ to enforce detailed balance. The second equation tracks the mediator, which is sourced by annihilation and depleted by cannibal reactions. The third equation expresses energy conservation within the hidden sector and determines the evolution of the hidden-sector temperature $T_h$. Under the Maxwell--Boltzmann statistics  adopted throughout this work, the energy density and pressure of each species are fully specified by its number density and temperature through $\rho_i = B_1^i\,n_i$ $\rho_i + P_i =B_2^i\, n_i$, Eqs.~\eqref{eq:MB_rho_P} and~\eqref{eq:B1}, so that Eq.~\eqref{eq:boltz_rho} can be rewritten as an evolution equation for the hidden-sector temperature $T_h$ given $n_X$ and $n_Y$. 

Finally, the equations above are written for self-conjugate $X$ (Majorana fermion), following the conventions of~\cite{Berlin:2016gtr}. When $X$ is not self-conjugate, i.e. a Dirac fermion, $n_X$ counts both particles and antiparticles and the cross sections $\sigmav_X$ and $\sigmavv_{YXX\to XX}$ are each replaced by half their value to avoid double-counting identical pairs in the initial state.


\subsubsection{Reformulation in terms of chemical potentials and $\xi$}
\label{sec:chempot_formulation}

Rather than evolving the number densities directly, it is advantageous to reformulate the system in terms of the dimensionless chemical potentials $\bar{\mu}_i \equiv \mu_i/\TD$, defined through
\begin{equation}
    Y_i = Y_i^{\text{eq}}\,e^{\bar{\mu}_i}\,, \qquad i \in \{X, Y\}\,,
    \label{eq:chempot_def}
\end{equation}
where $Y_i = n_i/s$ is the comoving yield, $s$ is the SM entropy density and $Y_i^{\text{eq}} = n^{\rm eq}_{{\rm MB},\,i}/s$ its equilibrium value. This parametrization has two key benefits. First, in the early universe when both species are in chemical equilibrium, $\bar{\mu}_i = 0$ exactly, avoiding the need to track the cancellation between $Y_i$ and $Y_i^{\text{eq}}$ --- a severe numerical challenge when both are large and nearly equal. Second, during freeze-out the dimensionless chemical potential grows linearly while the yield departs exponentially, so $\bar{\mu}_i$ is a much more slowly varying quantity and hence better suited as an ordinary differential equation (ODE) variable.

We take as independent variable $x \equiv m_X/T_{\rm SM}$ and promote
the dark-to-SM temperature ratio $\ln\xi$ to a dynamical degree of
freedom, obtaining a coupled system of three ODEs:
\begin{align}
    \frac{d\bar{\mu}_X}{dx} &= A + B\,\frac{d\ln\xi}{dx} \label{eq:dmuX} \\
    \frac{d\bar{\mu}_Y}{dx} &= C + D\,\frac{d\ln\xi}{dx} \label{eq:dmuY} \\
    \frac{d\ln\xi}{dx} &= \frac{E + F\,A + G\,C}{1 - F\,B - G\,D} \label{eq:dlnxi_full}
\end{align}
where the six coefficient functions $A$--$G$ depend on $(x,\,\xi,\,\bar\mu_X,\,
\bar\mu_Y)$ and are constructed from the thermally averaged cross
sections, the equilibrium thermodynamic quantities ($B_1$, $B_2$,
$\lambda$), and the SM entropy-conservation factor $\tilde g$ as defined in Secs.~\ref{sec:cosmology} and~\ref{sec:models}. The
explicit forms for these coefficients are given in Eqs.~\eqref{eq:Acoeff}--\eqref{eq:Gcoeff}
below. The full derivation of this system from the number-density Boltzmann equations~\eqref{eq:boltz_nX}--\eqref{eq:boltz_rho}, including the extension to nonzero mediator decay width, is given in Appendix~\ref{app:derivation}.
 
\paragraph{Dark-matter coefficients.}
%
\begin{align}
  A &= -\frac{s\,\tilde g}{H_{\rm tot}\,x}\,\langle\sigma v\rangle_{X} Y_X^{\rm eq}e^{\bar\mu_X}
       \left[1- e^{2(\bar\mu_Y-\bar\mu_X)}
       \right]
       + \frac{\lambda_X - 3\tilde g}{x}\,,
  \label{eq:Acoeff}
  \\
  B &= -\lambda_X\,.
  \label{eq:Bcoeff}
\end{align}
%
The first term in $A$ is the collision integral for $X\bar X \leftrightarrow YY$, written in a form that makes the approach to chemical equilibrium ($\bar\mu_X = \bar\mu_Y$) manifest. The second term encodes the redshifting of the equilibrium yield, with $\lambda_X \equiv d\ln n_X^{\rm eq}/d\ln T_h$, Eq.~\eqref{eq:lambda}, capturing the temperature sensitivity of the equilibrium density and $\tilde g$, Eq.~\eqref{eq:gtilde}, which accounts for the running of the SM entropy degrees of freedom.
 
\paragraph{Mediator coefficients.}
%
\begin{align}
  C &= \frac{\tilde g}{H_{\rm tot}\,x}\,\left\{
       s\,\langle\sigma v\rangle_{X}\, \frac{(Y^{\rm eq}_X)^2}{Y^{\rm eq}_Y}e^{2\bar\mu_X-\bar\mu_Y}\,\left[1-e^{2(\bar\mu_Y-\bar\mu_X)}\right]
       - \Gamma_{\rm can}\bigl(1 - e^{-\bar\mu_Y}\bigr)\right\}
       + \frac{\lambda_Y - 3\tilde g}{x}\,,
  \label{eq:Ccoeff}
  \\
  D &= -\lambda_Y\,,
  \label{eq:Dcoeff}
\end{align}
The first contribution in $C$ is the annihilation source, equal and opposite to the corresponding term in $A$. The second involves the total cannibal reaction rate $\Gamma_{\rm can}$ defined in Eq.~\eqref{eq:Gamma_can_std}, multiplied by the factor $(1 - e^{-\bar\mu_Y})$ which vanishes identically when $Y$ is in chemical equilibrium ($\bar\mu_Y = 0$) and acts as a restoring force that drives $\bar\mu_Y$ back toward zero whenever cannibal reactions are fast. The last term plays the same role as in $A$, accounting for the redshifting of $Y_Y^{\rm eq}$.
 
\paragraph{Temperature-equation coefficients.}
These arise from projecting the energy-conservation equation onto
$d\ln\xi/dx$:
%
\begin{align}
  E &= \frac{1}{x}
       \left(1 - 3\tilde g\,
             \frac{n_X\,B_2^X + n_Y\,B_2^Y}{\tilde D}
       \right),
  \label{eq:Ecoeff}
  \\
  F &= -\frac{B_1^X\,n_X}{\tilde D}\,,
  \qquad
  G = -\frac{B_1^Y\,n_Y}{\tilde D}\,,
  \label{eq:Gcoeff}
\end{align}
%
with the energy denominator,
\begin{equation}
  \tilde D
  = T_h\!\left(
      n_X\,\frac{dB_1^X}{dT_h}
    + n_Y\,\frac{dB_1^Y}{dT_h}
    \right)
  + B_1^X\,n_X\,\lambda_X
  + B_1^Y\,n_Y\,\lambda_Y\,.
  \label{eq:Dtilde}
\end{equation}
%
Here, the coefficient $E$ captures the adiabatic cooling of the hidden sector due to Hubble expansion, while $F$ and $G$ encode the feedback of changes in $n_X$ and $n_Y$ on the hidden-sector temperature. All three are controlled by the denominator $\tilde D$, which measures the total heat capacity of the dark plasma. The thermodynamic quantities $B_1$, $B_2$, and $\lambda$ used here are defined
in Sec.~\ref{sec:bessel}.

%======================================================================
\subsubsection{Solver Implementation}
\label{sec:solver}
%======================================================================

The system of Eqs.~\eqref{eq:dmuX}--\eqref{eq:dlnxi_full} is solved by the function \texttt{solve\_boltzmann\_chempot\_3phase}, which employs a three-phase strategy that exploits the natural hierarchy of timescales in the hidden sector. The DM is the first species to depart from chemical equilibrium, at a time when cannibal processes are still fast enough to hold the mediator at its equilibrium abundance. Only later do the cannibal reactions themselves become inefficient, and the mediator too falls out of equilibrium, entering the full freeze-out regime. Solving the full three-variable system from the outset would introduce unnecessary stiffness in the early evolution, where the mediator equation is effectively algebraic. The three-phase decomposition avoids this by progressively activating degrees of freedom as each process falls out of equilibrium.

In Phase~1, both species are in chemical equilibrium, i.e. $\bar\mu_X = \bar\mu_Y = 0$, and only the hidden-sector temperature evolves. The system reduces to a single ODE for $\ln\xi$, obtained from Eq.~\eqref{eq:dlnxi_full} with all equilibrium densities evaluated at $T_h$. In formulas:
\begin{equation}
    \frac{d\ln\xi}{dx}\bigg|_{\text{eq}} = \frac{1}{x}\left(1 - 3\tilde{g} \,\frac{n_Y^{\text{eq}}\,B_2^Y + n_X^{\text{eq}}\,B_2^X}{\TD\left(B_1^X\,\frac{dn_X^{\text{eq}}}{d\TD} + \frac{dB_1^X}{d\TD}\,n_X^{\text{eq}} + B_1^Y\,\frac{dn_Y^{\text{eq}}}{d\TD} + \frac{dB_1^Y}{d\TD}\,n_Y^{\text{eq}}\right)}\right)
    \label{eq:dlnxi_eq}
\end{equation}

The solver exits this phase when the dimensionless annihilation rate $\Gamma_{\rm ann}/H \equiv (s/H_{\rm tot})\,\langle\sigma v\rangle_X Y_X$ drops below a configurable threshold safely set at $5\times10^9$ by default (the parameter \texttt{Gamma\_over\_H\_ann} in the code).


In Phase~1.5, the DM has begun to depart from equilibrium but cannibal reactions remain fast enough to enforce $\bar\mu_Y = 0$. Setting $\bar\mu_Y = 0$
and $d\bar\mu_Y/dx = 0$ in Eqs.~\eqref{eq:dmuX}--\eqref{eq:dlnxi_full} eliminates the mediator equation entirely and simplifies the coefficients $A$ and $C$, since all factors of $Y_Y$ reduce to $Y_Y^{\rm eq}$.  The system then reduces to two ODEs for $(\ln\xi,\,\bar\mu_X)$, with the mediator abundance fixed at its equilibrium value. This intermediate phase is the key ingredient that resolves the stiffness near the annihilation freeze-out transition. It ends when $\Gamma_{\rm can}/H$ (see Eq.~\eqref{eq:Gamma_can_std}) drops below a configurable threshold, set to unity by default (the parameter \texttt{Gamma\_over\_H\_can} in the code), at which point the cannibal processes can no longer maintain the mediator in chemical equilibrium.


In Phase~2, all three variables $(\ln\xi,\,\bar\mu_X,\,\bar\mu_Y)$ evolve via the full system of Eqs.~\eqref{eq:dmuX}--\eqref{eq:dlnxi_full} with the coefficients given in Eqs.~\eqref{eq:Acoeff}--\eqref{eq:Dtilde}. For large mass ratios $r\equiv m_X/m_Y$, the mediator freezes out much later than the DM, and the integration range is extended accordingly to $x_{\rm max} = \max(10^5,\,r^2\times10^4)$.

The initial conditions are set at $x_0 = x_{\rm min}$ set by default to $10^{-3}$, deep in the equilibrium regime. The chemical potentials start at zero, $\bar\mu_X(x_0) = \bar\mu_Y(x_0) = 0$, with the initial hidden to SM temperature ratio determined by entropy conservation from a high reference temperature $T_{\rm ini}=10^{14}$~GeV by default. In formulas we get:
\begin{equation}
  \ln\xi(x_0) = \frac{1}{3}\ln\!\left(\frac{g_{*S}(m_X/x_0)}{g_{*S}(T_{\rm ini})}\right) + \ln\xi_{\rm ini}\,,
  \label{eq:xi_init}
\end{equation}
where $\xi_{\rm ini}$ is the corresponding asymptotic hidden to SM temperature ratio at the reference temperature, a free parameter of the model. It is important to emphasize that the precise value of $T_{\rm ini}$ is not relevant, provided it lies well above all mass scales in the problem and above the highest threshold at which $g_{\star S}$ changes, so that $g_{\star S}(T_{\rm ini})$ has reached its asymptotic value.

The integration is terminated once the system has converged. Convergence is checked for $x > 2\,x_{\rm FO}$, where $x_{\rm FO}$ is the value of x
at which $\bar\mu_X$ first exceeds $0.05$, by requiring that the hidden-sector temperature tracks the non-relativistic scaling expected when all species have frozen out, namely that:
\begin{equation}
  \frac{d\ln\xi}{d\ln x}\bigg|_{\rm NR} = 1 - 2\tilde g\,.
  \label{eq:dlnxi_NR}
\end{equation}
Specifically, the solver terminates once $|d\ln\xi/d\ln x - (1-2\tilde g)|$ falls below a configurable tolerance, set to $10^{-2}$ by default (the parameter \texttt{convergence\_threshold} in the code).

\subsection{Late-time background evolution and mediator decay}
\label{sec:background}

Once the Boltzmann system has converged, the comoving yields of $X$ and $Y$ are effectively frozen and the hidden-sector temperature has settled onto the non-relativistic scaling of Eq.~\eqref{eq:dlnxi_NR}. The subsequent evolution is governed by the decay of the mediator into SM particles, $Y\to\mathrm{SM}$, with a decay width $\Gamma_Y \propto \varepsilon^2$ whose explicit form depends on the portal model (see Sec.~\ref{sec:models}) and is now the only relevant timescale. This stage is handled by a separate solver, \texttt{solve\_background}, which evolves the system in the number of e-folds $N = \ln(a/a_0)$ using three log-scaled variables, namely $y_0 = \ln(T/T_0)$, $y_1 =\ln(\rho_Y/\rho_{Y,0})$, and $y_2 = \ln(\rho_X/\rho_{X,0})$, where the subscript "$0$" denotes values at the start of the background evolution. The evolution equations read
\begin{align}
    &\frac{d\ln T}{dN} = -\frac{1}{\gtilde}\left(1 - \frac{\rho_Y\,\Gamma_Y}{3H_{\rm tot}\,s\, T}\right)
    \label{eq:dlnTdN} \\
    &\frac{d\ln\rho_Y}{dN} = -\left(3 + \frac{\Gamma_Y}{H_{\rm tot}}\right)
    \label{eq:dlnrhoYdN} \\
    &\frac{d\ln\rho_X}{dN} = -3 - \frac{\sigmav_{\text{s-wave}}\,\rho_X}{\mX\,H_{\rm tot}}
    \label{eq:dlnrhoXdN}
\end{align}
with the total Hubble rate $H_{\rm tot}$ is evaluated at each step from Eq.~\eqref{eq:Htot}. The first equation describes the evolution of the SM temperature, which now receives an entropy injection term from the decay $Y\to\mathrm{SM}$ that slows the usual $T\propto a^{-1}$ redshifting. The second combines Hubble dilution of the mediator energy density with its depletion by decay, while the third tracks the DM energy density including a residual annihilation term that is typically negligible at this stage but is retained for completeness.

The key observable is the total entropy injected into the SM bath, quantified by the ratio $S_{\rm ratio} = S_{\rm final}/S_0$ of the final to initial comoving entropy $S = s\,a^3$. When the mediator decays promptly ($\Gamma_Y \gg H$), the entropy injection is negligible and $S_{\rm ratio} \approx 1$. Instead, for late decays $S_{\rm ratio}$ can be $\gg 1$ and the DM yield is diluted accordingly to $Y_X^{\rm final}=\rho_X^{\rm final}/(m_X s_{\rm final})  \approx Y_{X,0}/S_{\rm ratio}$. The evolution is terminated when $\rho_Y/\rho_{\rm tot} < 10^{-12}$, ensuring that the mediator has fully decayed and the entire entropy injection has been accounted for.

\paragraph{Decoupling the portal coupling from the Boltzmann evolution.}
A key design choice of the solver is to treat the mediator as stable during the Boltzmann evolution of Sec.~\ref{sec:solver}, neglecting the portal coupling entirely. This separation is motivated by efficiency. The Boltzmann system is numerically expensive to solve due to its stiffness and the wide dynamic range involved, so it is highly advantageous to solve it once for a given set of hidden-sector couplings and masses, and then scan over $\varepsilon$ using only the fast background evolution to determine the value that reproduces the observed DM relic abundance while satisfying BBN and direct-detection bounds. We have verified that this two-stage procedure agrees at the sub-percent level with a full joint evolution that includes $\Gamma_Y$ in the Boltzmann equations throughout.
 
For sufficiently large values of $\varepsilon$, however, the decay rate $\Gamma_Y$ can become comparable to the Hubble rate already during the Boltzmann evolution, meaning that some fraction of the mediator population should have decayed before the background solver begins. To account for this, the code provides an optional correction (\texttt{correct\_Y\_decay}), which retroactively steps through the Boltzmann trajectory and computes the cumulative decay and entropy injection that was missed. The surviving mediator fraction $f_{\rm survive}$ and the accumulated entropy correction $S_{\rm corr}$ are used to adjust the initial conditions for the background evolution according to
\begin{equation}
  Y_{Y,0} \to Y_{Y,0}\,f_{\rm survive}\,,\qquad
  T_0 \to T_0\left(\frac{S_{\rm corr}}{S_0}\right)^{1/3},\qquad
  Y_{X,0} \to \frac{Y_{X,0}}{S_{\rm corr}/S_0}\,.
  \label{eq:Ydecay_correction}
\end{equation}
This correction is negligible for most of the parameter space of interest in this work, namely for $\varepsilon \lesssim 10^{-10}$ where $\Gamma_Y \ll H$ throughout the Boltzmann evolution, but it becomes important for larger values of $\varepsilon$ where $\Gamma_Y/H \sim \mathcal{O}(1)$ near the end of the freeze-out phase.


%======================================================================
\section{Mapping the Portal Coupling: Predictions and Constraints}
\label{sec:mapping}
%======================================================================

The framework developed in Sec.~\ref{sec:evolution} provides, for any choice of the hidden-sector parameters --- masses $m_X$ and $m_Y$, and the hidden-sector couplings (e.g.\ $\alpha_X$ for the vector portals, or $\lambda_s,\,\lambda_p$ for the Higgs portal) --- the late-time DM yield $Y_X^{\rm final}$ and the frozen mediator abundance $Y_Y^{\rm frozen}$ as functions of the portal coupling~$\varepsilon$. Because $\varepsilon$ enters only through the mediator decay width, it controls the late-time entropy injection that dilutes the relic abundance but does not affect the hidden-sector freeze-out itself, which is governed entirely by the hidden-sector couplings. This separation makes $\varepsilon$ the natural variable through which to express both predictions and observational constraints.

In this section we describe how the toolkit uses this dependence to determine three quantities in the $(m_X,\,\varepsilon)$ plane, at fixed hidden-sector couplings and mass ratio $r = m_X/m_Y$:
\begin{itemize}
    \item the value $\varepsilon_{\rm relic}(m_X)$ that reproduces the observed DM relic abundance $\Omega_c h^2 = 0.12$ --- our \emph{prediction line} (Sec.~\ref{sec:eps_relic});
    \item the lower bound $\varepsilon_{\rm BBN}(m_X)$ from Big Bang Nucleosynthesis (BBN) hadronic injection by late mediator decays (Sec.~\ref{sec:bbn});
    \item the upper bound $\varepsilon_{\rm DD}(m_X)$ from spin-independent DM--nucleon scattering experiments (Sec.~\ref{sec:DD}).
\end{itemize}
The intersection of the relic line with the excluded regions determines the viable parameter space of the model. A separate requirement --- that the hidden sector remain thermally decoupled from the SM bath --- sets a thermalization floor on~$\varepsilon$ that we discuss in Sec.~\ref{sec:thermal_floor}.

\subsection{Dark Matter Relic Abundance}
\label{sec:eps_relic}

For a given choice of hidden-sector parameters, the value of the portal coupling that reproduces the observed DM relic abundance is determined by the function \texttt{find\_epsilon\_DMRelic}. The procedure runs the Boltzmann solver of Sec.~\ref{sec:solver} once to obtain the frozen-out yields, and then performs a one-dimensional root search on the background evolution of Sec.~\ref{sec:background} for the value of $\varepsilon$ that satisfies
\begin{equation}
  Y_X^{\rm final}(\varepsilon) = Y_{\rm target}\,, \qquad
  Y_{\rm target} \equiv \frac{\Omega_c h^2\,\rho_c}{s_0\,m_X}\,,
  \label{eq:Y_target}
\end{equation}
with $\Omega_c h^2 = 0.12$, $\rho_c = 1.05375\times 10^{-5}\,h^2\,\mathrm{GeV/cm^3}$, and $s_0 = 2891\,\mathrm{cm^{-3}}$. The root search is performed on $\log_{10}\varepsilon$ via Brent's method (\texttt{scipy.optimize.brentq}) over a default range $\log_{10}\varepsilon \in [-20,\,-6]$.

The physical interpretation of this procedure is straightforward. At fixed masses and hidden-sector couplings, the freeze-out fixes both $Y_X^{\rm frozen}$ and $Y_Y^{\rm frozen}$ at the end of the Boltzmann phase. The role of $\varepsilon$ is then to set the lifetime of $Y$ and hence the magnitude of the entropy injected into the visible sector once $Y$ decays. Larger $\varepsilon$ corresponds to a shorter lifetime, less entropy injection, and a less diluted relic. There is therefore a unique $\varepsilon = \varepsilon_{\rm relic}(m_X)$ for which the final yield matches $Y_{\rm target}$, provided $Y_X^{\rm frozen}$ is itself \emph{at least} as large as $Y_{\rm target}$ --- otherwise no amount of entropy dilution can recover the observed abundance.


\paragraph{Minimum DM mass and prediction-line.}

The above observation implies the existence of a minimum DM mass $m_X^\star$ --- determined by the hidden-sector couplings and the mass ratio $r = m_X/m_Y$ --- below which the model cannot reproduce $\Omega_c h^2 = 0.12$ regardless of how small $\varepsilon$ is taken. This critical mass is defined implicitly by
\begin{equation}
    m_X^\star\,Y_X^{\rm frozen}(m_X^\star) = \frac{\Omega_c h^2\,\rho_c}{s_0}\,.
    \label{eq:mXstar_def}
\end{equation}
For $m_X < m_X^\star$ the annihilation cross section is too large and the hidden-sector freeze-out alone underproduces the observed relic, so no value of $\varepsilon$ can recover it. For $m_X > m_X^\star$ the freeze-out overproduces, and a unique $\varepsilon$ is required to inject enough entropy to dilute the excess down to $Y_{\rm target}$; smaller $\varepsilon$ means a longer mediator lifetime, more entropy injection, and hence a larger DM mass that can be accommodated. The prediction line $\varepsilon_{\rm relic}(m_X)$ therefore curves away from $m_X^\star$ toward larger masses as $\varepsilon$ decreases.

At $m_X = m_X^\star$ itself, the prediction line enters a \emph{plateau regime}, namely the range of portal couplings large enough that the mediator decays promptly --- so that entropy dilution is negligible --- yet small enough that the decay width does not perturb the hidden-sector freeze-out. Within this window the final yield is independent of $\varepsilon$, and the prediction line is a vertical segment at $m_X = m_X^\star$. Below the lower edge of the plateau, $\varepsilon_{\rm min}$, the mediator lifetime becomes long enough for entropy injection to set in and the prediction line curves toward larger masses. For plotting purposes the vertical segment is extended upward to $\varepsilon \sim \mathcal{O}(1)$; the caveats associated with this extrapolation --- in particular the role of the thermalization floor --- are discussed at the end of this subsection.

The function \texttt{find\_mXstar} computes both $m_X^\star$ and $\varepsilon_{\rm min}$ for a given mass ratio~$r$. It first estimates the mass scale by inverting the standard freeze-out relation, $m_X^{\rm est}=\,\left[\sigmav_0\,(mY)_{\rm relic}\,\Mpl\right]^{1/2}$, where $\sigmav_0$ is the model's $s$-wave thermally averaged cross section evaluated at a reference DM mass of $1\,{\rm GeV}$ and $(mY)_{\rm relic} = \Omega_c h^2\,\rho_c/s_0$. The function then runs the full Boltzmann solver followed by the background evolution at \texttt{n\_refine} mass points (default $\texttt{n\_refine} = 5$) equally spaced in $\log_{10} m_X$ over the range $[0.1\,m_X^{\rm est},\; 5\,m_X^{\rm est}]$ as default, using a reference coupling $\varepsilon_{\rm ref} = 0.1$ together with the post-hoc decay correction. This correction (corresponding to \texttt{correct\_Y\_decay=True} in the code) retroactively accounts for $Y\to\mathrm{SM}$ decays along the Boltzmann trajectory and adjusts the handoff conditions for the background phase; it saturates to the plateau value for any $\varepsilon_{\rm ref}$ large enough that the mediator has fully decayed by the end of the integration, so the result is insensitive to the specific choice of reference coupling. A linear interpolation in $\left[\log_{10} m_X,\,\log_{10}(m_X\,Y_X^{\rm final})\right]$ then yields $m_X^\star$, and $\varepsilon_{\rm min}$ is obtained by a subsequent root search on the background entropy ratio, defined as the smallest portal coupling for which $S_{\rm ratio}(\varepsilon_{\rm min})/S_{\rm plateau} - 1 < \delta_S$, with $S_{\rm plateau} \equiv S_{\rm ratio}(\varepsilon_{\rm ref})$ and a default tolerance $\delta_S = 10^{-3}$.

\paragraph{Validity and caveats.}
Two systematic effects limit the accuracy of the plateau construction. Above the thermalization floor $\varepsilon_{\rm therm}$ (Sec.~\ref{sec:thermal_floor}), the hidden and visible sectors enter kinetic and chemical contact, so that $Y$ is replenished thermally and the hidden-sector freeze-out picture underlying the calculation of $m_X^\star$ no longer applies. In that regime the freeze-out proceeds as a standard WIMP in equilibrium with the SM bath, yielding a critical mass that is only slightly larger than the value computed here; the correction is minute and does not affect the exclusion picture. In addition, below the thermalization floor the plateau is also not perfectly flat. In fact, as $\varepsilon$ increases toward $\varepsilon_{\rm therm}$, the growing decay width $\Gamma_Y$ begins to perturb the hidden-sector freeze-out, introducing a percent-level drift in the predicted $m_X^\star$. This drift grows with the hidden-sector coupling, since a stronger coupling delays freeze-out to lower SM temperatures, where the entropy injected by the decay of the mediator has a proportionally larger impact on the relic abundance. We have verified, by comparison with a fully self-consistent joint evolution that includes the mediator decay throughout the Boltzmann phase (see Sec.~\ref{app:numerics}), that the resulting bias on $m_X^\star$ remains below the $10\%$ level even for the largest hidden-sector couplings explored in this work. Given the magnitude of these effects, the vertical-line representation of the prediction at $m_X = m_X^\star$ is an adequate approximation throughout the parameter space considered here.

\subsection{Constraints from Big Bang Nucleosynthesis}
\label{sec:bbn}

The second quantity of interest is the lower bound on the portal coupling set by BBN. In the hidden-sector scenario, the hidden-sector freeze-out leaves behind a sizeable comoving abundance of mediators $Y$ which, for sufficiently small $\varepsilon$, can survive long enough to decay during or after BBN and inject hadronic energy into the primordial plasma. If this injection is too large, it spoils the successful predictions of the light-element abundances, and the corresponding region of parameter space is excluded. The code adopts the constraints derived in~\cite{Kawasaki:2017bqm}, which are tabulated as an upper bound on the energy yield per decay $\zeta \equiv m_Y\,Y_Y$ (in GeV) as a function of the mediator lifetime $\tau_Y\equiv 1/\Gamma_Y$ (in s) for several reference masses and two representative hadronic injection channels, $u\bar u$ and $b\bar b$. A two-dimensional log--log interpolation (\texttt{scipy.interpolate.RegularGridInterpolator}) is built in the $(\log_{10}(m_Y/\mathrm{TeV}),\,\log_{10}\tau_Y)$ plane over six reference masses ($30\,\mathrm{GeV}$, $100\,\mathrm{GeV}$, $1\,\mathrm{TeV}$, $10\,\mathrm{TeV}$, $100\,\mathrm{TeV}$, $1000\,\mathrm{TeV}$). Outside the tabulated mass range the bound is extended via a power-law extrapolation:
\begin{equation}
    \zeta_{\max}(m_Y) = \zeta_{\max}(m_{\rm anchor})\,(m_Y/m_{\rm anchor})^\alpha
\end{equation}
where the slope $\alpha$ is extracted numerically as the local power-law index $\Delta\left(\log_{10}\zeta_{\max}\right)/\Delta\left(\log_{10}m_Y\right)$ between adjacent tabulated masses at fixed $\tau_Y$, evaluated at the edge of the grid closest to the extrapolation region.
For $m_Y > 1000\,\mathrm{TeV}$ (anchored to the highest tabulated point), both channels converge to $\alpha_{\rm high}\simeq 0.70$. This value follows from the QCD fragmentation multiplicity scaling $N_{\rm had}\sim m_Y^{0.3}$~\cite{Kawasaki:2017bqm}, which implies that a heavier particle produces more hadrons per decay, so fewer decays are needed to inject the same hadronic energy, and the upper bound on the total energy yield relaxes as $\zeta_{\max}\propto m_Y^{1-0.3}$. For $m_Y < 30\,\mathrm{GeV}$ (anchored to the lowest tabulated point), the slopes are reduced to $\alpha_{\rm low}=0.50$ for $u\bar u$ and $0.34$ for $b\bar b$. The departure from the asymptotic value of $0.7$ arises because at lower masses the energy per nucleon in the shower, $E_N\sim m_Y/N_{\rm had}\sim m_Y^{0.7}$, drops into a regime where the hadrodissociation cross-sections decrease with energy, so that each individual hadron becomes less effective at disrupting light elements and the bound tightens more slowly with mass than the multiplicity argument alone would predict. The particularly low $b\bar b$ slope reflects an additional suppression from the proximity to the $b$-quark threshold ($2m_b\approx 8.4\,\mathrm{GeV}$), which sharply reduces shower multiplicity relative to $u\bar u$.

Given the interpolated bound, for a given $(m_X,\,m_Y)$ and fixed hidden-sector couplings, the minimum portal coupling compatible with BBN, $\varepsilon_{\rm BBN}(m_X)$, is obtained by a one-dimensional root search on $\log_{10}\varepsilon$ of the condition
\begin{equation}
    m_Y\,Y_Y^{\rm frozen} = \zeta_{\max}\!\left(m_Y,\,\tau_Y(\varepsilon)\right)\,,
\end{equation}
with $Y_Y^{\rm frozen}$ read off from the Boltzmann solver output of Sec.~\ref{sec:solver}.

Across essentially the entire parameter space considered in this work, namely for $m_X\gtrsim m_X^\star$, the mediator abundance left over by the hidden-sector freeze-out is large enough that $m_Y\,Y_Y^{\rm frozen}$ exceeds $\zeta_{\max}(\tau_Y)$ for \emph{any} lifetime within the BBN-sensitive range. As a consequence, the BBN bound reduces to the simple requirement that $\tau_Y \lesssim 0.1$--$0.2\,\mathrm{s}$, above which charged pions from the mediator's hadronic shower survive long enough to drive $p\leftrightarrow n$ interconversion before thermalizing in the electromagnetic plasma~\cite{Kawasaki:2017bqm}. The BBN constraints associated with longer lifetimes, namely hadrodissociation of ${}^4\mathrm{He}$ ($\tau\sim 10^2$--$10^7\,\mathrm{s}$) and photodissociation ($\tau\gtrsim 10^4\,\mathrm{s}$), are never probed in our parameter space, as the large mediator abundance ensures that the constraint is already saturated at the minimum $p\leftrightarrow n$ interconversion threshold.


\paragraph{Validity and caveats.}
Two approximations enter the BBN bound as implemented. First, the power-law extrapolation below $30\,\mathrm{GeV}$ is conservative, as at lower masses the hadronic activity per decay is suppressed faster than the power law captures, so the true $\zeta_{\max}$ is larger (less constraining) than the extrapolated value and the code does not under-exclude parameter space. The extrapolation above $1000\,\mathrm{TeV}$ rests on the asymptotic QCD multiplicity scaling and is better motivated. Second, the bound is applied on each channel independently, weighted by the branching ratio, but leptonic channels carry no hadronic constraint in the short-lifetime regime relevant here, and the remaining hadronic final states (up-type and down-type quarks, gluons, electroweak bosons, $hh$) all produce similar pion-dominated showers and are grouped under the $u\bar u$ curves from~\cite{Kawasaki:2017bqm}, so that only the relative weight of $b\bar b$ against this $u\bar u$-like combination matters. In practice, both approximations have minimal impact on the final result, especially since the large mediator abundance characteristic of the hidden-sector scenarios under consideration cause the constraint to collapse to a single lifetime threshold $\tau_Y\lesssim 0.1$--$0.2\,\mathrm{s}$, which is insensitive to the detailed shape of the $\zeta_{\max}(\tau_Y)$ curve and hence to the interpolation and extrapolation choices above.

\subsection{Constraints from Direct Detection}
\label{sec:DD}


The third quantity of interest is the upper bound on the portal coupling from spin-independent DM--nucleon scattering experiments. At fixed hidden-sector parameters, the SI cross section $\sigma_{\rm SI}(\varepsilon)$ is a monotonic function of $\varepsilon$ through the portal-induced DM--nucleon amplitude discussed in Sec.~\ref{sec:models} for each model; its explicit form is Eq.~\eqref{eq:VP_sigma_SI} for the \texttt{VectorPortal}, Eq.~\eqref{eq:BL_sigma_SI} for the \texttt{BLPortal} and \texttt{BaryonPortal}, and Eq.~\eqref{eq:HP_sigma_SI} for the \texttt{HiggsPortal}. The \texttt{LiLjPortal} has no tree-level coupling to nucleons and is therefore unconstrained. The code compares $\sigma_{\rm SI}(\varepsilon,\,m_X)$ against the 90\% CL upper limit $\sigma_{\rm limit}(m_X)$ from a chosen experiment, with three built-in options: the latest constraints reported by the LZ~\cite{LZ:2024zvo} and XENON~\cite{XENON:2025vwd} collaborations, and a pre-tabulated coherent neutrino scattering floor~\cite{OHare:2021utq} used for the ultimate reach estimate. Each limit is loaded from the corresponding tabulated data file and evaluated through a log--log interpolation in $(m_X,\,\sigma_{\rm SI})$.

For a given $m_X$ and fixed hidden-sector parameters, the maximum portal coupling compatible with the chosen constraint, $\varepsilon_{\rm DD}(m_X)$, is obtained by a one-dimensional root search on $\log_{10}\varepsilon$ of the condition
\begin{equation}
    \sigma_{\rm SI}\!\left(\varepsilon_{\rm DD},\,m_X\right) = \sigma_{\rm limit}(m_X)\,.
\end{equation}
For DM masses outside the validity range of the chosen limit table, the function returns \texttt{NaN}; for models with vanishing tree-level cross section, i.e. the \texttt{LiLjPortal}, it returns $+\infty$, reflecting the absence of a meaningful direct-detection constraint.

%======================================================================


\subsection{Thermalization Floor}
\label{sec:thermal_floor}

The hidden-sector freeze-out picture underlying this work assumes that the hidden and visible sectors are not in thermal equilibrium at the time of DM freeze-out. For completeness, the code provides an estimate of the portal coupling above which this assumption breaks down, following the approach of Ref.~\cite{Hu:2025thq}. The condition is that the mediator decay rate into SM final states exceeds the Hubble rate at the freeze-out temperature $T_{\rm FO}\equiv m_X/x_{\rm FO}$ (with a default $x_{\rm FO}=20$),
\begin{equation}
    \Gamma(Y\to\mathrm{SM}) \geq H(T_{\rm FO})\,,
\end{equation}
which, since all partial widths scale as $\varepsilon^2$, yields the thermalization floor
\begin{equation}
    \varepsilon_{\rm therm} = \sqrt{\frac{H(T_{\rm FO})}{\Gamma_{\rm norm}(Y\to\mathrm{SM})}}\,,
    \label{eq:eps_therm}
\end{equation}
where $\Gamma_{\rm norm}$ is the total mediator decay width to SM particles evaluated at $\varepsilon=1$. For portal couplings above $\varepsilon_{\rm therm}$ the two sectors equilibrate, so that while DM still annihilates predominantly via $XX\to YY$, both species now share a common temperature with the SM bath and the hidden-sector dynamics that give rise to entropy dilution no longer applies.

\appendix

\section{Additional approximations and Validity}
\label{sec:assumptions}
 
This appendix collects the two main additional assumptions that are common to the entire Boltzmann system and that have not been discussed in the main text, namely the use of Maxwell–Boltzmann (MB) statistics for all hidden-sector species, and the assumption of internal kinetic equilibrium at a single hidden-sector temperature $T_h$. We then summarize the dominant sources of numerical error and delineate the region of parameter space over which the code has been validated.
 
\paragraph{Maxwell--Boltzmann statistics.}
All equilibrium distributions, collision integrals, and the chemical-potential parametrisation $Y_i = Y_i^{\rm eq}\,e^{\bar\mu_i}$, Eq.~\eqref{eq:chempot_def}, rely on MB statistics. The deviations from the full quantum distributions are largest in the relativistic regime, $x_h\equiv m/T_h \lesssim \mathcal{O}(1)$, where both species are still in chemical equilibrium, and become exponentially small by the time freeze-out occurs for $x_h\gtrsim 20-30$. As a whole, the MB contributes $\lesssim 1\%$ uncertainty on the computed DM relic abundance $\Omega_{X} h^2$, which is well below other systematics.
 
\paragraph{Kinetic equilibrium within the hidden sector.}
All hidden-sector species are assumed to share a common temperature $\TD$ at all times, so that the phase-space distribution of each species is fully characterised by its mass, chemical potential, and $\TD$. This requires elastic scattering processes within the hidden sector, such as $XY \to XY$, $XX \to XX$, $YY \to YY$, to remain fast compared to the Hubble rate throughout the evolution. In the models considered here the elastic rates scale as $\alphaX^2$ and are parametrically larger than the $2\to2$ annihilation and $3\to2$ cannibal rates that govern the chemical evolution, so kinetic equilibrium is maintained to a good approximation. The assumption would break down in scenarios with very small hidden-sector couplings or large mass hierarchies, where the elastic scattering rate could drop below $H$ before chemical freeze-out is complete. No explicit check of this condition is performed by the code; it is the user's responsibility to verify that the elastic rates remain sufficiently fast for a particular choice of hidden-sector parameters.
 
\paragraph{Numerical precision and regions of reduced accuracy.}
The code has been tested over the parameter ranges $\mX\sim5\GeV-100\TeV$, $\mY\sim 0.5-70\TeV$, mass ratios $r=\mX/\mY$ between
1.3 and 10, and hidden-sector couplings $\alphaX\gtrsim 10^{-4}$ (or $\lambda_{s,p} \gtrsim 10^{-2}$ for the Higgs portal). Within this region the relic abundance is accurate at the sub-percent level, dominated by the ODE integration error. Outside this range, the most common source of difficulty is either (i) the very late freeze-out which occurs for small hidden-sector couplings ($\alphaX \lesssim 10^{-4}$) or (ii) the prolonged cannibal equilibrium of the mediator at large mass ratios ($r \gtrsim 10$), which delays its departure from equilibrium over an extended range in $x$. In both cases the convergence criterion of Eq.~\eqref{eq:dlnxi_NR} may not be reached within the default integration range, requiring an increased $x_{\rm max}$ or finer grids. Together with the model-specific uncertainties discussed in the main text, the dominant sources of numerical error and their estimated magnitudes are collected in Table~\ref{tab:numerics}.

\begin{table}
\centering
\begin{tabular}{ll}
\toprule
Source & Estimated error \\
\midrule
ODE integration (Radau, \texttt{rtol}=$10^{-8}$) & $\lesssim 0.01\%$ on $Y_X$ \\
MB statistics vs.\ full integral ($z \sim 1$--$3$) & $1$--$10\%$ on $\neqMB$; $\lesssim 1\%$ on $\Omega_X h^2$ \\
$\gstar$ interpolation & $\lesssim 1\%$ (smooth tables) \\
\texttt{HDECAY} decay widths & $\lesssim 5\%$ (see Sec.~\ref{sec:HP_SM}) \\
BBN bound interpolation & $\lesssim 10\%$ (grid density dependent) \\
$\Delta_{1,2,3}$ cannibal coefficients & $\mathcal{O}(1)$ uncertainty \\
\bottomrule
\end{tabular}
\caption{Dominant sources of numerical uncertainty in the relic abundance calculation. All estimates refer to the validated parameter region described in the text.}
\label{tab:numerics}
\end{table}
 


\section{Details on the numerical implementation and alternative solvers}
\label{app:numerics}

This appendix collects additional technical details on the numerical integration of the Boltzmann system presented in Sec.~\ref{sec:boltzmann}. We first describe three alternative solver strategies that are retained in the codebase for testing and cross-validation purposes but are not used in production. We then summarize the numerical settings common to all solvers, including the integration grid, tolerances, and regularization procedures.

\subsection{Alternative solvers}
\label{app:legacy}

Three additional solvers are implemented alongside the primary three-phase solver described in Sec.~\ref{sec:solver}. They differ in how aggressively the phase structure is exploited and in the choice of dynamical variables during the intermediate epoch.

The two-phase solver (\texttt{solve\_boltzmann\_chempot\_2phase}) skips the intermediate Phase~1.5 entirely, transitioning directly from full equilibrium (Phase~1) to the complete three-variable system (Phase~2) when $\Gamma_{\rm ann}/H$ drops below $10^7$. This introduces unnecessary stiffness near the annihilation freeze-out transition, where $Y$ is still effectively in equilibrium but $X$ is departing, precisely the difficulty that the intermediate phase of the primary solver is designed to resolve.

The hybrid solver (\texttt{solve\_boltzmann\_hybrid}) does include an intermediate phase, but formulates it in terms of the physical yields $Y_X$ and $Y_Y$ rather than the dimensionless chemical potentials:
\begin{align}
\frac{dY_X}{dx} &= -\frac{s\,\tilde g}{H_{\rm tot}\,x}\,\sigmav_X\left(Y_X^2 - (Y_X^{\rm eq})^2\,\frac{Y_Y^2}{(Y_Y^{\rm eq})^2}\right), \label{eq:dYX_hybrid} \\
\frac{dY_Y}{dx} &= -\frac{dY_X}{dx} - \frac{\tilde g}{H_{\rm tot}\,x}\,\Gamma_{\rm can}\,(Y_Y - Y_Y^{\rm eq})\,,
\label{eq:dYY_hybrid}
\end{align}
where $\Gamma_{\rm can}$ is the cannibal rate defined in Eq.~\eqref{eq:Gamma_can_std}.

It transitions to the chemical-potential variables of Phase~2 when either $\Gamma_{\rm can}/H$ falls below a threshold or the fractional departure $|Y_X - Y_X^{\rm eq}|/Y_X^{\rm eq}$ exceeds a configurable value. While conceptually straightforward, tracking $Y_Y$ directly can suffer from numerical loss of precision when it spans many orders of magnitude during the cannibal epoch.

The $Y$-equilibrium solver (\texttt{solve\_boltzmann\_Yeq}) is the simplest of the four, assuming that the mediator always tracks its equilibrium abundance throughout the entire evolution. It solves only two variables, $(\ln\xi,\,\bar\mu_X)$, using equations identical to Phase~1.5 of the primary solver. By construction, this solver cannot capture the cannibal freeze-out of $Y$ or its subsequent departure from equilibrium, which affects both the late-time temperature evolution and the mediator abundance at the onset of the background evolution stage.

\subsection{Numerical settings}
\label{app:num_details}

All solvers use \texttt{scipy.integrate.solve\_ivp} with the implicit Runge--Kutta method Radau (L-stable, suitable for stiff problems). The maximum step size is set to $1$ in $x$, with relative tolerance $10^{-8}$ and absolute tolerance $10^{-15}$ during Phase~2 (relaxed to $10^{-12}$ during Phase~1). Rather than integrating from $x_{\min}$ to $x_{\max}$ in a single call, the evolution is broken into segments over a predefined grid of breakpoints,
\begin{equation}
    x \in \{10^{-2},\, 10^{-1},\, 0.99\} \cup \mathrm{linspace}(1,\,10,\,20) \cup \{10.5,\, 11.5,\, \ldots,\, 99.5\} \cup \{200,\, 500,\, \ldots,\, x_{\rm max}\}\,,
\end{equation}
with $500$ evaluation points per segment. This segmented approach improves the robustness of the solver. The integrator restarts fresh at each breakpoint, which helps it adapt its internal step size to the rapidly changing stiffness of the system. Phase-switching conditions are nevertheless checked continuously within each segment.


Several clipping and regularisation procedures are applied to prevent numerical overflow and underflow. The temperature ratio is restricted to $\ln\xi \in [-20,\,20]$ and the dimensionles chemical potentials to $|\bar\mu_i| \leq 10/x$ per Hubble time. Exponentials of yields are clipped to the range $[-700,\,700]$, and densities that fall below $10^{-300}$ are replaced with $\pm 10^{-300}$, preserving the original sign.

In addition to the default convergence criterion of the primary solver described in Sec.~\ref{sec:solver}, an alternative mode is available in which the solver instead monitors the rate of change of the DM yield and terminates when $|d\ln Y_X/d\ln x|$ falls below the configurable tolerance (the parameter \texttt{convergence\_threshold} in the code). This mode can be selected via the \texttt{convergence\_mode} option.

%======================================================================
%======================================================================
\section{Derivation of the Boltzmann system in chemical-potential variables}
\label{app:derivation}
%======================================================================

This appendix provides the step-by-step derivation of the coupled ODE system in Eqs.~\eqref{eq:dmuX}--\eqref{eq:dlnxi_full} and of the coefficient functions~\eqref{eq:Acoeff}--\eqref{eq:Gcoeff} summarised in Sec.~\ref{sec:boltzmann}. We retain the mediator decay width $\Gamma_Y$ throughout for completeness, and because the joint evolution that includes it is used for the cross-check discussed in Sec.~\ref{app:numerics}. As shown in Sec.~\ref{sec:background}, setting $\Gamma_Y = 0$ is an excellent approximation in the regime of interest ($\varepsilon \lesssim 10^{-10}$, where $\Gamma_Y \ll H$ throughout the freeze-out process), and doing so at any point in the expressions below recovers the system solved by the default Boltzmann integrator.
 
Starting from the dimensionless chemical-potential parametrisation $Y_i = Y_i^{\rm eq}\,e^{\bar\mu_i}$ of Eq.~\eqref{eq:chempot_def} and differentiating with respect to $x = m_X/T_{\rm SM}$,
\begin{equation}
  \frac{dY_i}{dx} = Y_i\!\left[\frac{d\bar\mu_i}{dx} + \frac{d\ln Y_i^{\rm eq}}{dx}\right]\,.
\end{equation}
The log-derivative of the equilibrium yield follows from $Y_i^{\rm eq} = n^{\rm eq}_{{\rm MB},\,i}(\TD)/s(T_{\rm SM})$, the relation $\ln\TD = \ln\xi + \ln m_X - \ln x$ (so that $d\ln\TD/dx = d\ln\xi/dx - 1/x$), and SM entropy conservation $d\ln s/dx = -3\tilde g/x$, with $\tilde g$ defined in Eq.~\eqref{eq:gtilde},
\begin{equation}
  \frac{d\ln Y_i^{\rm eq}}{dx}
  = \lambda_i\!\left(\frac{d\ln\xi}{dx} - \frac{1}{x}\right) + \frac{3\tilde g}{x}\,,
  \label{eq:dlnYeq_app}
\end{equation}
where $\lambda_i$ is given by Eq.~\eqref{eq:lambda}. Combining the two,
\begin{equation}
  \frac{dY_i}{dx} = Y_i\!\left[\frac{d\bar\mu_i}{dx} + \lambda_i\!\left(\frac{d\ln\xi}{dx} - \frac{1}{x}\right) + \frac{3\tilde g}{x}\right],
  \label{eq:dYi_chempot}
\end{equation}
which is the identity used repeatedly below to trade $dY_i/dx$ for the variables $d\bar\mu_i/dx$ and $d\ln\xi/dx$.
 
The number-density Boltzmann equations~\eqref{eq:boltz_nX}--\eqref{eq:boltz_nY}, supplemented by the mediator decay $Y\to\mathrm{SM}$ as a number sink $-\Gamma_Y\,n_Y$ in the $Y$ equation, are converted to comoving yields via $\dot Y_i = (Hx/\tilde g)\,dY_i/dx$ and $\dot n_Y = s\dot Y_Y - 3Hs\,Y_Y$, giving
\begin{align}
  \frac{dY_X}{dx}
    &= -\frac{s\,\tilde g}{H_{\rm tot}\,x}\,\sigmav_X\left(Y_X^2 - \frac{(Y_X^{\rm eq})^2}{(Y_Y^{\rm eq})^2}\,Y_Y^2\right), \label{eq:BEQ_YX_app}\\[4pt]
  \frac{dY_Y}{dx}
    &= -\frac{dY_X}{dx} - \frac{\tilde g}{H_{\rm tot}\,x}\,\Gamma_{\rm can}\,(1 - e^{-\bar\mu_Y}) - \frac{\tilde g}{H_{\rm tot}\,x}\,\Gamma_Y\,Y_Y\,.
  \label{eq:BEQ_YY_app}
\end{align}
The last term in~\eqref{eq:BEQ_YY_app} is the decay contribution, arising from the $-\Gamma_Y\,n_Y$ sink rewritten as $\dot Y_Y\big|_{\rm decay} = -\Gamma_Y\,Y_Y$ and converted to the $x$-derivative. The hidden-sector energy conservation equation~\eqref{eq:boltz_rho} also picks up a source term, since each decay removes the average energy $B_1^Y$ from the hidden plasma and deposits it into the SM bath,
\begin{equation}
  \dot\rho_h + 3H(\rho_h + P_h) = -\Gamma_Y\,B_1^Y\,n_Y\,,
  \label{eq:energy_cons_app}
\end{equation}
where we have used $\rho_i = B_1^i\,n_i$ from Eq.~\eqref{eq:MB_rho_P}.
 
\paragraph{Dark-matter and mediator equations.}
 
The DM equation is unaffected by $\Gamma_Y$. Dividing Eq.~\eqref{eq:BEQ_YX_app} by $Y_X$, applying the identity~\eqref{eq:dYi_chempot} on the left-hand side, and expressing the collision integral in chemical potentials via $Y_i = Y_i^{\rm eq}\,e^{\bar\mu_i}$, one obtains
\begin{equation}
  \frac{d\bar\mu_X}{dx} + \lambda_X\!\left(\frac{d\ln\xi}{dx} - \frac{1}{x}\right) + \frac{3\tilde g}{x}
  = -\Gamma_{\rm ann}\left(e^{\bar\mu_X} - e^{2\bar\mu_Y - \bar\mu_X}\right),
\end{equation}
with $\Gamma_{\rm ann} \equiv (s\,\tilde g)/(H_{\rm tot}\,x)\,\sigmav_X\,Y_X^{\rm eq}$. Rearranging in the form $d\bar\mu_X/dx = A + B\,d\ln\xi/dx$ directly reproduces Eqs.~\eqref{eq:Acoeff}--\eqref{eq:Bcoeff} of the main text, with $A = -\Gamma_{\rm ann}(e^{\bar\mu_X} - e^{2\bar\mu_Y - \bar\mu_X}) + (\lambda_X - 3\tilde g)/x$ and $B = -\lambda_X$.
 
The mediator equation $d\bar\mu_Y/dx = (dY_Y/dx)/Y_Y - d\ln Y_Y^{\rm eq}/dx$ decomposes into three collision contributions plus the redshifting of $Y_Y^{\rm eq}$. The annihilation source, obtained from the $-dY_X/dx$ term in Eq.~\eqref{eq:BEQ_YY_app}, is
\begin{equation}
  \frac{1}{Y_Y}\!\left(-\frac{dY_X}{dx}\right) = \Gamma_{\rm ann}\,\frac{Y_X^{\rm eq}}{Y_Y^{\rm eq}}\,e^{-\bar\mu_Y}\!\left(e^{2\bar\mu_X} - e^{2\bar\mu_Y}\right).
\end{equation}
The $3\to 2$ cannibal contribution yields the standard $-(\tilde g/x)(\Gamma_{\rm can}/H_{\rm tot})(1 - e^{-\bar\mu_Y})$ term already appearing in Eq.~\eqref{eq:Ccoeff}, and the decay term in~\eqref{eq:BEQ_YY_app} gives simply $-(\tilde g/x)(\Gamma_Y/H_{\rm tot})$, independent of $\bar\mu_Y$ and of $d\ln\xi/dx$. Subtracting $d\ln Y_Y^{\rm eq}/dx$ via~\eqref{eq:dlnYeq_app} and collecting terms gives $d\bar\mu_Y/dx = C + D\,d\ln\xi/dx$, with
\begin{align}
  C &= \Gamma_{\rm ann}\,\frac{Y_X^{\rm eq}}{Y_Y^{\rm eq}}\,e^{-\bar\mu_Y}\!\left(e^{2\bar\mu_X} - e^{2\bar\mu_Y}\right) - \frac{\tilde g}{H_{\rm tot}\,x}\,\Gamma_{\rm can}(1 - e^{-\bar\mu_Y}) - \frac{\tilde g}{H_{\rm tot}\,x}\,\Gamma_Y + \frac{\lambda_Y - 3\tilde g}{x}\,, \label{eq:C_app}\\[4pt]
  D &= -\lambda_Y\,.
\end{align}
Relative to the $\Gamma_Y = 0$ system of Eq.~\eqref{eq:Ccoeff}, the coefficient $C$ acquires a constant decay drain $-(\tilde g/x)(\Gamma_Y/H_{\rm tot})$ that depletes the mediator yield below its chemical-equilibrium value.
 
\paragraph{Temperature equation.}
 
The temperature equation is obtained from the hidden-sector energy conservation law~\eqref{eq:energy_cons_app}. Writing $\rho_h = \sum_i s\,B_1^i(\TD)\,Y_i$, the time derivative is
\begin{equation}
  \dot\rho_h = \dot s\sum_i B_1^i\,Y_i + s\sum_i\!\left[\frac{dB_1^i}{d\TD}\,\dot\TD\,Y_i + B_1^i\,\dot Y_i\right].
\end{equation}
Substituting this into Eq.~\eqref{eq:energy_cons_app}, using $\dot s = -3Hs$, $3H(\rho_h + P_h) = 3Hs\sum_i B_2^i\,Y_i$, and the identity $B_2^i = B_1^i + \TD$ from Eq.~\eqref{eq:B1}, the $B_1^i$ terms on the two sides cancel. Converting all cosmic-time derivatives to $x$-derivatives via $\dot f = (Hx/\tilde g)\,df/dx$ and dividing through by $sHx/\tilde g$, one finds
\begin{equation}
  \frac{d\TD}{dx}\sum_i \frac{dB_1^i}{d\TD}\,Y_i
  + \sum_i B_1^i\,\frac{dY_i}{dx}
  + \frac{3\tilde g\,\TD}{x}\sum_i Y_i
  = -\frac{\tilde g}{H_{\rm tot}\,x}\,\Gamma_Y\,B_1^Y\,n_Y\,,
  \label{eq:energy_x_app}
\end{equation}
where we have used $H\approx H_{\rm tot}$ in the Hubble-dilution piece on the left, an excellent approximation since the hidden sector is at most a subdominant component of the total energy density at the epoch of interest, while keeping the full $H_{\rm tot}$ in the decay term on the right.
 
Replacing $dY_i/dx$ with the identity~\eqref{eq:dYi_chempot} and using $d\TD/dx = \TD(d\ln\xi/dx - 1/x)$, the terms proportional to $(d\ln\xi/dx - 1/x)$ combine into the energy denominator $\tilde D$ already defined in Eq.~\eqref{eq:Dtilde}, the terms involving $d\bar\mu_i/dx$ give $\sum_i B_1^i\,n_i\,d\bar\mu_i/dx$, and the remaining $1/x$ terms combine into $(3\tilde g/x)\sum_i B_2^i\,n_i$ after using $B_2^i = B_1^i + \TD$. Equation~\eqref{eq:energy_x_app} then becomes
\begin{equation}
  \tilde D\!\left(\frac{d\ln\xi}{dx} - \frac{1}{x}\right)
  + \sum_i B_1^i\,n_i\,\frac{d\bar\mu_i}{dx}
  + \frac{3\tilde g}{x}\sum_i B_2^i\,n_i
  = -\frac{\tilde g}{H_{\rm tot}\,x}\,\Gamma_Y\,B_1^Y\,n_Y\,.
  \label{eq:energy_compact_app}
\end{equation}
Solving for $d\ln\xi/dx$ yields the form $d\ln\xi/dx = E + F\,d\bar\mu_X/dx + G\,d\bar\mu_Y/dx$ with
\begin{align}
  E &= \frac{1}{x}\!\left(1 - 3\tilde g\,\frac{n_X\,B_2^X + n_Y\,B_2^Y}{\tilde D}\right) - \frac{\tilde g}{H_{\rm tot}\,x}\,\Gamma_Y\,\frac{B_1^Y\,n_Y}{\tilde D}\,, \label{eq:E_app}\\[4pt]
  F &= -\frac{B_1^X\,n_X}{\tilde D}\,, \qquad G = -\frac{B_1^Y\,n_Y}{\tilde D}\,,
\end{align}
which reproduces the coefficients of Eqs.~\eqref{eq:Ecoeff}--\eqref{eq:Gcoeff} with the additional decay contribution to $E$. Physically, this extra term accelerates the cooling of the hidden sector as decaying mediators carry energy out of the dark plasma while depositing entropy into the SM bath, driving $\TD$ down faster than Hubble expansion alone.
 
\paragraph{Decoupled system.}
 
Substituting $d\bar\mu_X/dx = A + B\,d\ln\xi/dx$ and $d\bar\mu_Y/dx = C + D\,d\ln\xi/dx$ into the expression for $d\ln\xi/dx$ yields the decoupled master equation~\eqref{eq:dlnxi_full}, with the chemical-potential equations then following by back-substitution through~\eqref{eq:dmuX}--\eqref{eq:dmuY}. The algebraic structure of this decoupling is identical to the $\Gamma_Y = 0$ case; all decay dependence is absorbed into the modified coefficients $C$ and $E$ given in Eqs.~\eqref{eq:C_app} and~\eqref{eq:E_app}. Specifically, $C$ acquires a constant drain $-(\tilde g/x)(\Gamma_Y/H_{\rm tot})$ that depletes the mediator yield, and $E$ acquires an energy-loss term $-(\tilde g/x)(\Gamma_Y/H_{\rm tot})(B_1^Y n_Y/\tilde D)$ that accelerates the cooling. The coefficients $A$, $B$, $D$, $F$, $G$ and the energy denominator $\tilde D$ are unchanged, and in the limit $\Gamma_Y \ll H_{\rm tot}$ both corrections vanish, recovering the system solved by the default Boltzmann integrator.


%======================================================================
\section{Physical Constants}
\label{sec:constants}
%======================================================================
Table~\ref{tab:constants} collects the numerical values of all physical constants used in the code, taken from the Particle Data Group~\cite{PhysRevD.110.030001}.

\begin{longtable}{lll}
\toprule
\textbf{Symbol} & \textbf{Value} & \textbf{Description} \\
\midrule
\endhead
\multicolumn{3}{l}{\textit{Fundamental constants}} \\
$\Mpl$ & $2.435\times 10^{18}\GeV$ & Reduced Planck mass \\
$v_H$ & $246.22\GeV$ & Higgs vacuum expectation value \\
$m_N$ & $0.9389\GeV$ & Nucleon mass \\
$\hbar^2 c^2$ & $0.3894\times 10^{-27}\GeV^2\cm^2$ & Conversion factor \\
$\hbar$ & $6.582\times 10^{-25}\GeV\cdot\text{s}$ & Reduced Planck constant \\
\midrule
\multicolumn{3}{l}{\textit{Electroweak parameters}} \\
$\sin^2\theta_W$ & 0.23121 & Weinberg angle \\
$\alpha_{\text{EM}}$ & $1/137.0$ & Fine-structure constant \\
$e$ & 0.3029 & EM coupling \\
$g_W$ & 0.6520 & Weak coupling \\
$g_Y^{\text{SM}}$ & 0.3577 & Hypercharge coupling \\
$\alpha_s$ & 0.1179 & Strong coupling at $M_Z$ \\
\midrule
\multicolumn{3}{l}{\textit{Boson masses}} \\
$m_h$ & $125.10\GeV$ & Higgs mass \\
$M_Z$ & $91.1876\GeV$ & $Z$ boson mass \\
$M_W$ & $80.379\GeV$ & $W$ boson mass \\
$\Gamma_Z$ & $2.4952\GeV$ & $Z$ total width \\
\midrule
\multicolumn{3}{l}{\textit{Fermion masses}} \\
$m_e$ & $0.511\times 10^{-3}\GeV$ & Electron \\
$m_\mu$ & $0.1057\GeV$ & Muon \\
$m_\tau$ & $1.777\GeV$ & Tau \\
$m_u$ & $2.16\times 10^{-3}\GeV$ & Up quark \\
$m_d$ & $4.67\times 10^{-3}\GeV$ & Down quark \\
$m_s$ & $93.4\times 10^{-3}\GeV$ & Strange quark \\
$m_c$ & $1.27\GeV$ & Charm quark \\
$m_b$ & $4.18\GeV$ & Bottom quark \\
$m_t$ & $172.69\GeV$ & Top quark \\
\midrule
\multicolumn{3}{l}{\textit{Cosmological observables}} \\
$s_0$ & $2891\cm^{-3}$ & Entropy density today \\
$\rho_c$ & $1.054\times 10^{-5}\GeV\,h^2\cm^{-3}$ & Critical density \\
$\Och$ & 0.12 & Observed DM relic density \\
$m_Y^{\text{relic}}$ & $\Och\,\rho_c/s_0$ & Reference relic yield $\times$ mass \\
\bottomrule
\caption{Physical constants and Standard Model parameters used throughout the code.}
\label{tab:constants}
\end{longtable}


%======================================================================
\section{Code Architecture}
\label{sec:architecture}
%======================================================================
The code is organized around a strict separation between physics and numerics. Model instances (\texttt{VectorPortal}, \texttt{HiggsPortal}, etc.) are stateless, carrying only their physical parameters and providing rate methods with no reference to a \texttt{Cosmology} instance. The \texttt{BoltzmannSolver} combines a \texttt{Cosmology} and a \texttt{DarkSectorModel} instance, reading all model-specific physics through the abstract interface. Mass scans are handled via a factory pattern in which \texttt{find\_mXstar} takes a callable that creates model instances for different mass points. Heavy initializations such as \texttt{HDECAY} compilation and BBN data loading are deferred until first use and cached thereafter. The full module layout is shown below.

\begin{verbatim}
source/
  __init__.py              Public API exports
  constants.py             Physical constants, fermion tables
  cosmology.py             Cosmology class (g*, Hubble, thermo)
  model.py                 DarkSectorModel ABC + 5 portals
  solver.py                BoltzmannSolver + find_mXstar
  bbn.py                   BBN hadronic injection bounds
  direct_detection.py      LZ and neutrino floor interpolators
  SecludedDM.py            Backward-compatible re-exports

  phi_decay/
    __init__.py            ensure_hdecay_ready, public API
    phi_decay_hybrid.py    Three-regime stitching
    hdecay_interface.py    HDECAY Fortran wrapper
    scalar_portal/         Low-mass regime (< 5 GeV)
      api/model.py         LightScalar / HeavyScalar
      decay/               Decay channel implementations
      data/                Form factors, QCD, constants

  data/
    direct_detection_data/ LZ, neutrino floor tables
    bbn_constraints_decay/ 12 BBN data files

  _legacy/
    DM_Relic_old.py        Deprecated monolithic version
\end{verbatim}


\section*{AI disclosure}
The initial draft of this document was generated by Claude (Anthropic) from an analysis of the codebase. The final version has been extensively rewritten and edited by the author; very little of the original machine-generated text remains.
\clearpage
\phantomsection
\addcontentsline{toc}{section}{References}
\bibliographystyle{utphys}
\bibliography{bibl}

\end{document}
